\documentclass[11pt,letterpaper]{article}
\usepackage[T1]{fontenc}
\usepackage[utf8]{inputenc}
\usepackage{amsmath,amsthm,mathtools}
\usepackage{newpxtext,newpxmath}
\usepackage[margin=1in,headheight=15pt]{geometry}
\usepackage{microtype}
\usepackage{booktabs,longtable,array}
\usepackage{xcolor}
\usepackage{graphicx}
\usepackage{float}
\usepackage{enumitem}
\usepackage{listings}
\usepackage{fancyhdr}
\usepackage{titlesec}
\usepackage[most]{tcolorbox}
\usepackage{hyperref}
\usepackage{bookmark}
\definecolor{Navy}{HTML}{18334B}
\definecolor{Teal}{HTML}{27656B}
\definecolor{Pale}{HTML}{F1F5F7}
\definecolor{CodeGray}{HTML}{525B63}
\hypersetup{colorlinks=true,linkcolor=Navy,citecolor=Teal,urlcolor=Teal,
 pdftitle={The Six-State Shuffle Bound: Three Independent Certificate Computations},
 pdfauthor={Research manuscript prepared for Vladimir Reshetnikov},
 pdfsubject={Regular languages, shuffle state complexity, computer-assisted proof}}
\urlstyle{same}
\titleformat{\section}{\Large\bfseries\color{Navy}}{\thesection}{0.65em}{}
\titleformat{\subsection}{\large\bfseries\color{Navy}}{\thesubsection}{0.65em}{}
\titleformat{\subsubsection}{\normalsize\bfseries\color{Navy}}{\thesubsubsection}{0.65em}{}
\pagestyle{fancy}
\fancyhf{}
\fancyhead[L]{\small\color{CodeGray}The six-state shuffle bound}
\fancyhead[R]{\small\color{CodeGray}Research draft}
\fancyfoot[C]{\small\thepage}
\renewcommand{\headrulewidth}{0.3pt}
\setlength{\parskip}{0.35em}
\setlength{\parindent}{1.15em}
\setlist{itemsep=2pt,topsep=4pt}
\setcounter{tocdepth}{2}
\numberwithin{equation}{section}
\newtheorem{theorem}{Theorem}[section]
\newtheorem{lemma}[theorem]{Lemma}
\newtheorem{proposition}[theorem]{Proposition}
\newtheorem{corollary}[theorem]{Corollary}
\newtheorem{conjecture}[theorem]{Conjecture}
\theoremstyle{definition}
\newtheorem{definition}[theorem]{Definition}
\newtheorem{example}[theorem]{Example}
\newtheorem{question}[theorem]{Question}
\theoremstyle{remark}
\newtheorem{remark}[theorem]{Remark}
\newcommand{\shuf}{\mathbin{\operatorname{shuf}}}
\newcommand{\scf}{\operatorname{sc}_{\mathrm{shuf}}}
\newcommand{\id}{\operatorname{id}}
\newcommand{\suppR}{\operatorname{rows}}
\newcommand{\suppC}{\operatorname{cols}}
\newcommand{\Hard}{\mathsf{Hard}}
\newcommand{\Valid}{\mathsf{Valid}}
\newcommand{\Span}{\mathsf{Span}}
\newcommand{\Cert}{\mathsf{Cert}}
\newcommand{\bits}{\operatorname{bits}}
\newcommand{\pop}{\operatorname{popcount}}
\newcommand{\Pow}{\mathcal{P}}
\newcommand{\code}[1]{\texttt{\detokenize{#1}}}
\newcommand{\Sym}{\operatorname{Sym}}
\newcommand{\Resid}{\mathsf{Residual}}
\newcommand{\PreCert}{\mathsf{PreCert}}
\newcommand{\dbA}{\textsf{A}}
\newcommand{\dbB}{\textsf{B}}
\newcommand{\dbC}{\textsf{C}}
\newtcolorbox{scopebox}{colback=Pale,colframe=Teal,boxrule=0.6pt,
 arc=1mm,left=3mm,right=3mm,top=2mm,bottom=2mm}
% Marks material that is deliberately kept alongside a parallel account of the
% same statement; the argument of each box says what that variant buys.
\newtcolorbox{variantbox}[1]{colback=Pale,colframe=Navy,boxrule=0.5pt,
 arc=1mm,left=3mm,right=3mm,top=2mm,bottom=2mm,
 before upper={\textbf{\color{Navy}Parallel account \##1.}\ }}
\lstset{basicstyle=\ttfamily\footnotesize,keywordstyle=\bfseries\color{Navy},
 commentstyle=\itshape\color{CodeGray},stringstyle=\color{Teal},
 frame=single,rulecolor=\color{black!20},backgroundcolor=\color{Pale},
 breaklines=true,columns=fullflexible,keepspaces=true,showstringspaces=false,
 tabsize=4,aboveskip=9pt,belowskip=9pt}
\makeatletter
\let\inputrows\@@input
\makeatother
\begin{document}
\hypersetup{pageanchor=false}
\begin{titlepage}
\thispagestyle{empty}
{\color{Teal}\sffamily\bfseries FORMAL LANGUAGES \quad / \quad COMPUTER-ASSISTED MATHEMATICS\par}
\vspace{0.7in}
{\Huge\bfseries\color{Navy}The Six-State\par Shuffle Bound\par}
\vspace{0.22in}
{\Large Three independent certificate computations,\par
one finite-to-infinite theorem,\par constructive reachability\par}
\vspace{0.4in}
{\large Research manuscript prepared for Vladimir Reshetnikov\par}
\vspace{0.1in}
{20 September 2026\par}
\vspace{0.4in}
\begin{abstract}
The shuffle operation interleaves two words without changing the order of the
letters within either word. Its exact worst-case deterministic state complexity
is the subject of a published conjecture. This manuscript gives a
computer-assisted proof of the following conjecture in the range
$\min(m,n)\le6$.

\begin{conjecture}[Published shuffle state-complexity conjecture]
\label{conj:main}
For all $m,n\ge2$, the worst-case deterministic state complexity of the shuffle
is
\begin{equation}\label{eq:F}
 F(m,n)=2^{mn-1}+2^{(m-1)(n-1)}(2^{m-1}-1)(2^{n-1}-1).
\end{equation}
\end{conjecture}

\noindent
This article proves that the bound \eqref{eq:F} is attained whenever $m,n\ge2$
and $\min(m,n)\le6$. The general conjecture is
not settled here. The shared core theorem is proved once, from a single
finite-to-infinite induction over the lexicographic rank $(m+n,|S|)$; the
finite premise of that induction is then discharged three separate times, by
three certificate databases that are genuinely independent of one another.
Database~\dbA{} contains 10,642 spanning predecessor witnesses found by an SMT
search and audited by an orbit-closure coverage enumeration. Database~\dbB{}
contains 10,619 witnesses on a provably minimal, degree-sorted set of
representatives; its target keys are a strict subset of \dbA{}'s, and not one
of the 10,619 shared keys carries the same witness, so each database
re-derives the other's finite lemma from scratch. Database~\dbC{} contains
453 witnesses on a strictly smaller irreducible class obtained by adding a
fifth, analytic reduction --- a balanced-triple contraction --- which shrinks
the finite obligation from widths at most twenty to widths at most eleven and
replaces the solver by a deterministic permutation enumeration. The
manuscript also proves three increasingly sharp access-word bounds,
$|S|+\min(m,n)-2$, the support form $E+\min(r,c)-2$, and the dual bound
$\min\{2(|S|-1),|S|+\min(m,n)-2\}$, and gives four distinguishability
arguments, one of them over a fixed three-letter alphabet. The archive
contains all three certificate databases, four independent checkers, four
reaching-word constructors, replayable examples, and audit reports.
\end{abstract}
\vfill
\begin{scopebox}
\textbf{Status and scope.} This is a computer-assisted research draft, not a
peer-reviewed or proof-assistant-verified manuscript. No proof-assistant
formalization and no independent human review is claimed for any part of it.
The six-state result extends the range explicitly established in the primary
sources located for this project; a targeted literature search found no prior
resolution of this strip, but that negative search result is not a proof of
novelty, and bibliographic priority remains provisional. No result for all
$m,n\ge7$, or for a fixed small alphabet, is claimed.
\end{scopebox}
\end{titlepage}
\hypersetup{pageanchor=true}
\tableofcontents
\newpage

\section{The problem and the result}
\label{sec:introduction}

\subsection{A concrete question about interleaving languages}
Let $\kappa(L)$ denote the number of states in the minimal \emph{complete}
deterministic finite automaton recognizing a regular language $L$. For two
words $u$ and $v$, their shuffle is the set of all interleavings that preserve
the internal order of each word. For instance,
\[
 ab\shuf c=\{abc,acb,cab\}.
\]
The shuffle of two languages is obtained by taking the union over all choices
of a word from each language. We write
\[
 \scf(m,n)=\max\{\kappa(K\shuf L):\kappa(K)=m,\ \kappa(L)=n\},
\]
where the maximum ranges over pairs of languages over a common finite
alphabet; the alphabet is allowed to depend on $m$ and $n$.

The chosen open problem is the conjecture that, for every $m,n\ge2$,
\begin{equation}\label{eq:conjectured-bound}
 \scf(m,n)=F(m,n):=
 2^{mn-1}+2^{(m-1)(n-1)}(2^{m-1}-1)(2^{n-1}-1).
\end{equation}
Brzozowski, Jir\'askov\'a, Liu, Rajasekaran, and Szyku\l a established the
upper bound and its attainment for $2\le m\le5$ with arbitrary $n\ge2$, as
well as for $m=n=6$ \cite{BJLRS2016}. Their reduction to reachable subsets of
a rectangular state space is the starting point here. Caron, Luque, and
Patrou subsequently developed a partition-based reformulation and an
antichain reduction, without settling the unrestricted conjecture
\cite{CLP2020}. These are the specific primary-source results against which
this manuscript is compared.

\subsection{What is established here}
\begin{theorem}[Six-state strip, computer-assisted]\label{thm:main}
For all integers $m,n\ge2$ with $\min(m,n)\le6$,
\[
 \scf(m,n)=F(m,n).
\]
In particular, for every $n\ge2$,
\begin{equation}\label{eq:six-simplified}
 \scf(6,n)=63\cdot2^{6n-6}-31\cdot2^{5n-5}.
\end{equation}
\end{theorem}

The part of this strip beyond the range explicitly obtained in the cited
work is $m=6$, $n\ge7$, and its transpose. The present proof does not consist
of checking a few additional rectangles. It reduces \emph{all} widths to a
finite collection of irreducible matrices and verifies a stronger,
root-independent predecessor property for every member of that collection.

A quantitative reachability result accompanies the state-count theorem.
For the universal alphabet introduced in Section~\ref{sec:automata}, every
valid subset $S$ has a reaching word of length at most
\begin{equation}\label{eq:word-bound-intro}
 |S|+\min(m,n)-2.
\end{equation}
Thus an exponentially large reachable subset automaton can nevertheless
have short access words. The bound is not asserted to be optimal, and
Section~\ref{sec:induction} in fact proves two strictly sharper statements
alongside it: a \emph{support form} $E+\min(r,c)-2$, in which $E=|S|$ and
$r,c$ count the nonempty rows and columns rather than the grid dimensions,
and a \emph{dual bound} $\min\{2(|S|-1),\,|S|+\min(m,n)-2\}$ whose first
branch is independent of the dimensions altogether.

The finite premise of the induction is discharged three times over, by three
certificate databases that are not reducible to one another
(Section~\ref{sec:computation}). Two of them classify the same irreducible
remainder by two different symmetry-reduction correctness proofs and share no
witness at all; the third classifies a strictly smaller remainder, obtained by
an extra analytic reduction proved in Section~\ref{sec:triple}. Each database
alone suffices for Theorem~\ref{thm:main}. Keeping all three is what makes the
finite step independently re-derivable rather than a single computation to be
taken on trust.

\subsection{Why this is a suitable finite attack}
Two aspects make this problem particularly amenable to a certificate-based
approach. First, a putative last letter is simply a pair of functions on
finite sets. Second, the difficult matrices have incomparable columns.
For six rows, an antichain contains at most $\binom63=20$ columns. This is a
bound on the width of a \emph{minimal obstruction}, not a restriction on the
widths covered by the theorem.

The central methodological choice is to search for a predecessor that meets
\emph{every} row and every column. Such a predecessor is valid regardless of
which row and column are distinguished as initial. This permits an unrooted
classification, while retaining a sound rooted reachability induction.
Without this precaution, quotienting by all row and column permutations can
silently lose the initial-state condition.

\subsection{Provenance of this manuscript}
\label{sec:provenance}
This document is a merge of four independently prepared research packages on
the same conjecture. Every one of them proved the shared core theorem stated
above; they are retained here because \emph{how} each proved it differs, and
those differences are mathematical content, not presentation.

\begin{itemize}
\item \code{shuffle-six-state-certificates} is the spine of the present text.
It proved the theorem from database~\dbA{}: 10,642 spanning predecessor
certificates over orbit-closure-checked representatives, discovered by an SMT
search, validated by a Python local checker that exports a verified-target
list, and audited for coverage by a separate C++ enumeration consuming only
that list. It contributed the manuscript's structure, the tagged-multigraph
reformulation, the Spanning Predecessor Conjecture, the certificate schema
and minimal verification core, and the review checklist.
\item \code{shuffle-six-state-strip} proved the theorem from the \emph{same}
database \dbA{} by the same proof: its certificate file and the one shipped
here agree in all 10,642 keys and all 10,642 witnesses, a fact recomputed
during this merge. It is retained for what it adds around that shared core:
the sharper support-form word bound with its explicit potential, a second
plain-integer serialization with a record-by-record agreement test, a
single-binary C++ verifier reporting totals across one through six rows, an
iterative explicit-stack constructor, a 111,984-word regression suite, a run
under undefined-behavior sanitization, and a standalone literature audit.
\item \code{shuffle-six-state-theorem} proved the theorem from
database~\dbB{}: an \emph{independent} finite computation over a different
canonicalization (degree-sorted representatives, proved to be exactly one per
orbit), yielding 10,619 witnesses of which not one coincides with \dbA{}'s
witness for the same target. It also proved an additional structural property
of its certificates and an entirely different distinguishability theorem over
a fixed three-letter alphabet.
\item \code{shuffle-six-state-complexity} proved the theorem by a materially
different proof: a fifth analytic reduction removes an entire class of
obstructions before any table is consulted, shrinking the finite obligation
to database~\dbC{} --- 453 witnesses over widths at most eleven --- and
replacing the solver by a deterministic, solver-free permutation enumeration
whose output was regenerated byte for byte. It also proves the dual word-length
bound and a fourth distinguishability argument.
\end{itemize}

Throughout, material carried over from a package other than the spine is
integrated into the spine's notation and numbering. Where two packages prove
the same statement by genuinely different machinery, both proofs are kept and
marked with a \emph{parallel account} box saying what the variant buys.

\subsection{Research status and novelty boundary}
The source review was conducted on 20 September 2026. The accessible 2016
paper, the 2019 preprint of the later journal article, and the journal's
publication record were checked. The mathematical source review used the
arXiv preprint of \cite{CLP2020} rather than the access-restricted journal
full text. Within \cite{BJLRS2016} the located positions are the abstract;
Table~1 on page~6 of the preprint, which marks the $6\times7$ and larger
six-row cases as unresolved; Lemmas~8--9 on pages~7--8; Theorem~10 on page~8;
and Lemma~12 with Corollary~13 on pages~9--10. In \cite{CLP2020} the term
``dense'' states refers to incomparable row and column supports, not to graph
density. Targeted searches did not identify a later resolution of the
six-state strip or of the full conjecture; those searches were limited by the
query formulations used and by indexing coverage, and did not reach
unpublished or unindexed work.

This is evidence for choosing the problem, not a guarantee of an exhaustive
literature search. Accordingly, the theorem is a mathematical claim supported
by the proof and artifacts below, while its claim to bibliographic novelty
remains provisional. The elementary reductions are presented with proofs
rather than with claims of individual priority. The manuscript should receive
independent mathematical and implementation review before being treated as an
established research publication.

\section{The universal shuffle automaton}
\label{sec:automata}

\subsection{From automata and grammars to a Boolean matrix}
For a positive integer $m$, put $[m]=\{0,1,\ldots,m-1\}$. Let the state sets
of two complete deterministic automata be $R=[m]$ and $C=[n]$. On a letter
$a$, suppose their transitions are $f_a:R\to R$ and $g_a:C\to C$.
The usual nondeterministic automaton for shuffle has state set $R\times C$
and transitions
\begin{equation}\label{eq:cell-transition}
 (i,j)\xrightarrow{a}\{(f_a(i),j),(i,g_a(j))\}.
\end{equation}
At each step, the current letter is assigned to exactly one of the two
component words. A run records these assignments. Consequently a run
reaching a pair of final states is equivalent to a decomposition of the
input into an accepted word of the first automaton and an accepted word of
the second, interleaved in their original orders. This proves that
\eqref{eq:cell-transition} recognizes shuffle.

There is a direct grammar interpretation. A complete deterministic automaton
is presented as a deterministic right-linear grammar: introduce one
nonterminal $A_q$ for each state $q$, add a production
$A_q\to aA_{\delta(q,a)}$ for each transition, and add $A_q\to\varepsilon$ at
each accepting state; the start nonterminal corresponds to the initial state.
Writing $X_i$ for the nonterminal of the first operand's state $i$, its
productions are $X_i\to aX_{f_a(i)}$ and, at final states, $X_i\to\varepsilon$.
In the shuffle grammar, nonterminals $Z_{i,j}$ have the two productions
\[
 Z_{i,j}\to aZ_{f_a(i),j},\qquad
 Z_{i,j}\to aZ_{i,g_a(j)},
\]
and an empty production at each pair of final states. Shuffle therefore
combines two descriptions in a different way from concatenation: at each
input letter one must decide which component produced it, and determinization
tracks all component-state pairs compatible with those choices. The present
question asks how large a deterministic finite-state description of this
interleaving language may have to be. It therefore lies squarely in
formal-language and regular-grammar theory, even though its remaining
obstruction is combinatorial.

The measure being bounded below is \emph{deterministic} state complexity. The
theorem is not a lower bound of the same size on nondeterministic automata or
on regular grammars, which can be much smaller, and it says nothing about the
time needed to recognize a particular input word.

The determinized automaton has subsets $S\subseteq R\times C$ as states.
We identify $S$ with its $m$-by-$n$ zero-one incidence matrix. A letter
$(f,g)$ acts by
\begin{equation}\label{eq:Delta}
 \Delta_{f,g}(S)
 =\{(f(i),j):(i,j)\in S\}\ \cup\
   \{(i,g(j)):(i,j)\in S\}.
\end{equation}
The initial subset is a singleton $\{(r,c)\}$, where $r$ and $c$ are the
initial states of the operands.

\subsection{The full transformation-pair alphabet}
Define
\begin{equation}\label{eq:alphabet}
 \Gamma_{m,n}=R^R\times C^C.
\end{equation}
Each pair of total functions is a distinct letter. The alphabet is finite,
with $m^m n^n$ letters. The two component automata read the same letter
$(f,g)$ by applying $f$ and $g$, respectively.

Any pair of complete $m$- and $n$-state automata uses some selection of these
transformation pairs, possibly with repeated pairs under different letter
names. Thus allowing all pairs cannot remove any reachable subset.
Conversely, the universal alphabet itself gives an admissible finite
witness alphabet. Its use is therefore legitimate for a worst-case
state-complexity theorem with no fixed alphabet bound. This universal model
is also used in \cite[Section~2.1]{BJLRS2016}.

In what follows, a ``word'' is a finite sequence of transformation pairs.
Letters are not charged according to the size of their function tables;
word length has its usual automata-theoretic meaning.

\subsection{Validity and the upper bound}
For $S\subseteq R\times C$, set
\[
 \suppR(S)=\{i:\exists j\ (i,j)\in S\},\qquad
 \suppC(S)=\{j:\exists i\ (i,j)\in S\}.
\]
A subset is \emph{valid relative to $(r,c)$} when
\[
 \Valid_{r,c}(S)\quad\Longleftrightarrow\quad
 r\in\suppR(S)\ \text{and}\ c\in\suppC(S).
\]
\begin{lemma}[Support persistence, in equality form]\label{lem:persistence}
For every subset $P$ and every letter $(f,g)$,
\begin{equation}\label{eq:support-equalities}
 \suppR(\Delta_{f,g}(P))=\suppR(P)\cup f(\suppR(P)),\qquad
 \suppC(\Delta_{f,g}(P))=\suppC(P)\cup g(\suppC(P)).
\end{equation}
In particular
$\suppR(P)\subseteq\suppR(\Delta_{f,g}(P))$ and
$\suppC(P)\subseteq\suppC(\Delta_{f,g}(P))$, so every subset reachable from
$\{(r,c)\}$ is valid relative to $(r,c)$.
\end{lemma}
\begin{proof}
The image cell $(i,g(j))$ contributes the first coordinate $i$, so it
contributes exactly $\suppR(P)$; the image cell $(f(i),j)$ contributes exactly
$f(\suppR(P))$. Their union is the first equality, and the second is the same
calculation with the coordinates exchanged. Inclusion follows, and applying it
successively to the letters of a reaching word, starting from the singleton,
gives validity.
\end{proof}

\begin{corollary}[Supports of a predecessor]\label{cor:pred-support}
If $\Delta_{f,g}(P)=S$, then $\suppR(P)\subseteq\suppR(S)$ and
$\suppC(P)\subseteq\suppC(S)$.
\end{corollary}
\begin{proof}
Immediate from \eqref{eq:support-equalities}.
\end{proof}

The equality form of the lemma is worth stating separately, because the
sharper word-length bound of Theorem~\ref{thm:support-bound} is an induction
on support counts and needs to know that supports never shrink along a
reduction, not merely that they never shrink along a reaching word.

Validity does not require $(r,c)\in S$, and the failure is not a corner case.
On a $2\times2$ grid, the letter whose two components are both the
transposition sends $\{(0,0)\}$ to $\{(1,0),(0,1)\}$, which meets row $0$ and
column $0$ without containing $(0,0)$. Losing sight of this distinction can
invalidate an induction that silently changes its initial point.

There is a second, execution-level proof that reachable implies valid, which
is worth recording because it exposes where completeness of the operand
automata is used. One possible nondeterministic execution gives every input
letter to the column operand, leaving the row coordinate equal to $r$
throughout; another gives every letter to the row operand, leaving the column
coordinate equal to $c$. Both executions are available only because the
operand transition functions are total.

\begin{proposition}[Counting valid subsets]\label{prop:upper}
There are exactly $F(m,n)$ valid subsets relative to a specified root.
Consequently $\scf(m,n)\le F(m,n)$.
\end{proposition}
\begin{proof}
Rename the root as $(0,0)$. Counting by inclusion and exclusion, the subsets
missing row~$0$ number $2^{(m-1)n}$, those missing column~$0$ number
$2^{m(n-1)}$, and those missing both number $2^{(m-1)(n-1)}$, so the valid
subsets number
\begin{equation}\label{eq:inclusion-exclusion}
 2^{mn}-2^{(m-1)n}-2^{m(n-1)}+2^{(m-1)(n-1)}.
\end{equation}
Equivalently, and this is the form used later: subsets containing $(0,0)$
contribute $2^{mn-1}$; a valid subset that omits $(0,0)$ must contain a
nonempty subset of the other $n-1$ cells of row~$0$, and a nonempty subset of
the other $m-1$ cells of column~$0$, while all $(m-1)(n-1)$ remaining cells
are optional. The two disjoint cases give exactly
\eqref{eq:conjectured-bound}, and factoring $2^{(m-1)(n-1)}$ out of
\eqref{eq:inclusion-exclusion} gives the same expression.
Lemma~\ref{lem:persistence} bounds the number of reachable determinized
states, and minimizing the deterministic automaton cannot increase that
number.
\end{proof}

This rederives the upper bound of \cite{BJLRS2016}. The actual work is to
show that validity is sufficient for reachability in the claimed range.

\section{Distinguishability requires no computation}
\label{sec:distinguishability}
Reachability is not by itself a state-complexity lower bound: two reachable
subset states might accept exactly the same suffixes. This section proves
distinguishability four times. Three of the four proofs are over the full
transformation-pair alphabet and differ in their machinery; the fourth is a
different theorem over a \emph{fixed three-letter} alphabet, and yields a
strictly stronger corollary. None of the four uses any computation, and none
depends on any later section. Throughout, assume $m,n\ge2$.

For the three universal-alphabet arguments, choose one final state $a\in R$
and one final state $b\in C$, so that a subset is accepting precisely when it
contains $(a,b)$.

\subsection{A forward two-letter construction}
\begin{lemma}[A two-letter cell test]\label{lem:celltest}
For every cell $(p,q)\in R\times C$, there is a word of length two accepted
from the singleton $\{(p,q)\}$ and from no other singleton.
\end{lemma}
\begin{proof}
Choose $a_0\in R\setminus\{a\}$, $b_0\in C\setminus\{b\}$, and
$c_0\in C\setminus\{q\}$. Let the second letter be $(h,k)$, where
\[
 h(i)=a_0\quad\text{for every }i,\qquad
 k(j)=\begin{cases}b,&j=q,\\ b_0,&j\ne q.\end{cases}
\]
Because $h$ never reaches $a$, a single source cell can produce the final
cell $(a,b)$ under this letter only through its second branch. That branch
reaches $(a,b)$ exactly from $(a,q)$.

Let the first letter be $(f,g)$, where
\[
 f(i)=\begin{cases}a,&i=p,\\ a_0,&i\ne p,\end{cases}
 \qquad g(j)=c_0\quad\text{for every }j.
\]
The second branch of this letter never reaches column $q$. Its first
branch reaches $(a,q)$ exactly from $(p,q)$. Therefore the two-letter word
$(f,g)(h,k)$ is accepted from precisely the desired source cell.
\end{proof}

\begin{corollary}[All different subsets are distinguishable]\label{cor:distinct}
Every two distinct subsets of $R\times C$ are distinguishable in the
universal shuffle automaton with final cell $(a,b)$.
\end{corollary}
\begin{proof}
Choose a cell in the symmetric difference of the two subsets. The word
from Lemma~\ref{lem:celltest} is accepted from a subset exactly when the
subset contains that cell, because the transition of a union is the union
of the transitions.
\end{proof}

\subsection{The same statement by a cellwise existential predecessor}
\begin{variantbox}{1}
The next proof establishes Lemma~\ref{lem:celltest} again, but it constructs
the two letters \emph{backwards} from the accepting cell, using an explicit
preimage computation instead of a forward case analysis. What it buys is a
closed formula, $\Delta_A^{-1}[\{(a,b)\}]=\{(p,b)\}$, that identifies the
whole intermediate state rather than only its relevant member; that formula
is reused in the single-cell separation tests of
Section~\ref{sec:algorithm}. It also isolates exactly where $m,n\ge2$ is
needed --- one must be able to choose an image avoiding a prescribed
coordinate.
\end{variantbox}

\begin{proof}[Second proof of Lemma~\ref{lem:celltest}]
For a letter $A=(f_A,g_A)$ write $\Delta_A^{-1}[T]$ for the
\emph{cellwise existential predecessor} of $T$: the set of individual cells
having at least one $A$-successor in $T$. It is not an inverse function on
subset states.

Choose $f_A$ with $f_A^{-1}(\{a\})=\{p\}$ and choose $g_A$ whose image avoids
$b$; both are possible because each coordinate set has at least two elements.
Then
\[
 \Delta_A^{-1}[\{(a,b)\}]
 =\bigl(f_A^{-1}(\{a\})\times\{b\}\bigr)
  \cup\bigl(\{a\}\times g_A^{-1}(\{b\})\bigr)
 =\{(p,b)\}.
\]
For a second letter $B=(f_B,g_B)$, choose $f_B$ whose image avoids $p$ and
choose $g_B$ with $g_B^{-1}(\{b\})=\{q\}$. Then
$\Delta_B^{-1}[\{(p,b)\}]=\{(p,q)\}$. Hence the word $BA$, read $B$ first, is
accepted from exactly the singleton $\{(p,q)\}$; by the union identity
\begin{equation}\label{eq:union}
 \Delta_{f,g}(X\cup Y)=\Delta_{f,g}(X)\cup\Delta_{f,g}(Y),
\end{equation}
which holds for a letter and hence for a word, a
subset accepts it exactly when it contains that cell.
\end{proof}

\begin{remark}
The two letters depend on the tested cell. Neither proof says that one fixed
pair of letters works for all tests, or that a two-letter alphabet attains
the state-complexity bound.
\end{remark}

\subsection{The same statement by enumerating run allocations}
\begin{variantbox}{2}
The third proof of the two-letter test reasons about the \emph{shuffle
semantics} directly instead of about subset transitions: it enumerates the
four ways of allocating the two letters between the two operands and
eliminates three of them. What it buys is that it never manipulates the
determinized transition at all, so it is the argument to check if one
suspects an error in \eqref{eq:Delta} itself; it also exhibits the separating
letters as an explicit table.
\end{variantbox}

Fix the final states $a=m-1$ and $b=n-1$, and use state $0$ as a nonfinal
state in each coordinate; these are distinct because $m,n\ge2$. Fix a cell
$(p,q)$ and choose any $d\in C\setminus\{q\}$. Define two letters by
\[
\begin{aligned}
 \alpha(i)&=\begin{cases}a,&i=p,\\0,&i\ne p,\end{cases}
 &\gamma(j)&=d, &\quad A&=(\alpha,\gamma),\\
 \zeta(i)&=0,
 &\beta(j)&=\begin{cases}b,&j=q,\\0,&j\ne q,\end{cases}
 &\quad B&=(\zeta,\beta).
\end{aligned}
\]

\begin{proof}[Third proof of Lemma~\ref{lem:celltest}]
Consider a singleton input $(i,j)$ and the four ways to allocate the letters
$A$ then $B$ to the row or the column operand. If $B$ goes to the row operand,
the final row coordinate is $0$, so acceptance is impossible; this eliminates
two allocations. If both letters go to the column operand, $A$ first puts the
column at $d\ne q$ and $B$ then puts it at $0$, again rejecting. The only
remaining allocation gives $A$ to the row operand and $B$ to the column
operand, and its final cell is $(\alpha(i),\beta(j))$, which equals $(a,b)$
exactly when $i=p$ and $j=q$. The statement for arbitrary subsets follows by
\eqref{eq:union}.
\end{proof}

\subsection{A fixed three-letter alphabet, and a stronger corollary}
\label{sec:three-letter}
\begin{variantbox}{3}
The following is not another proof of Lemma~\ref{lem:celltest}: it is a
different theorem. It uses a \emph{fixed} alphabet of three letters, whatever
$m$ and $n$ are, and it produces closed-form distinguishing words. What it
buys is twofold. First, it is the only distinguishability result here that
survives if one insists on a small alphabet, and it therefore reproduces the
construction of \cite{BJLRS2016} inside this manuscript instead of importing
it. Second, its conclusion is strictly stronger: \emph{all} $2^{mn}$ subsets
are pairwise distinguishable, reachable or not. Nothing in the other three
arguments gives that.
\end{variantbox}

Make $(m-1,n-1)$ the unique accepting cell and consider these three letters of
the universal alphabet:
\begin{center}
\begin{tabular}{lll}
\toprule
Letter & Row transformation & Column transformation\\
\midrule
$a$ & $i\mapsto i+1\pmod m$ & $j\mapsto0$\\
$b$ & $i\mapsto0$ & $j\mapsto j+1\pmod n$\\
$c$ & $0\mapsto1$, all $i\ge1\mapsto0$ & $j\mapsto n-1$\\
\bottomrule
\end{tabular}
\end{center}

\begin{lemma}[Singleton-distinguishing words]\label{lem:threeletter}
For every cell $(i,j)$ there is a word $w_{i,j}$ over $\{a,b,c\}$ accepted
from exactly that cell and from no other.
\end{lemma}
\begin{proof}
For a set of target cells, a letter's preimage consists of all sources with at
least one transition into it. Three uniqueness observations suffice. First,
when $r\ge1$ and $j\ge1$, the unique $a$-predecessor of $(r,j)$ is $(r-1,j)$,
because the column branch of $a$ always lands in column zero. Second, when
$r\ge1$ and $j\ge1$, the unique $b$-predecessor of $(r,j)$ is $(r,j-1)$,
because the row branch of $b$ always lands in row zero. Third, the unique
$c$-predecessor of $(1,0)$ is $(0,0)$, since $n-1\ne0$ and the only row sent
to row one is row zero.

Starting from the unique final cell and working backwards with these
observations gives
\begin{equation}\label{eq:distwords}
 w_{i,j}=\begin{cases}
 a^{m-1-i}b^{n-1-j},&j\ge1,\\[2pt]
 b\,a^{m-1-i}b^{n-2},&j=0,\ i\ge1,\\[2pt]
 c\,b\,a^{m-2}b^{n-2},&i=j=0.
 \end{cases}
\end{equation}
All exponents are nonnegative. In the first line the backward steps stay in
positive columns and the $b$ steps take place in the last row. In the second
line the extra initial $b$ moves the unique predecessor from column one to
column zero in a positive row. The last line uses the unique $c$-predecessor
of $(1,0)$.
\end{proof}

\begin{corollary}[All subsets, reachable or not]\label{cor:subsetsdistinct}
All $2^{mn}$ subsets of $R\times C$ are pairwise distinguishable, whether or
not they are reachable.
\end{corollary}
\begin{proof}
If two subsets differ, choose a cell in their symmetric difference. The word
$w$ of Lemma~\ref{lem:threeletter} is accepted from a subset exactly when the
subset contains that cell.
\end{proof}

\subsection{Minimality of the operands}
Both component automata over $\Gamma_{m,n}$ are themselves minimal.
Every component state is reachable using a constant transformation. Two
different states are distinguished by a transformation sending one to its
unique final state and the other to a nonfinal state. Thus their state
complexities really are $m$ and $n$, not merely bounded by these numbers.

It follows that proving reachability of all valid subsets in the universal
model proves equality in \eqref{eq:conjectured-bound}. The proof optimizes
the number of states, not the number of alphabet symbols: a separating
\emph{word} of length two does not mean that a two-symbol alphabet suffices
for all separations and all reaching words.

\section{Root-sensitive reductions}
\label{sec:reductions}
This section develops the reductions used by the induction. Containment
and degree-one reductions are part of the prior approach
\cite[Lemmas~8--9]{BJLRS2016}; complete versions are given here so that root
positions and the exact number of seed letters are explicit.

\subsection{Relabeling and restriction to a subgrid}
Let $\pi:R\to R'$ and $\sigma:C\to C'$ be bijections, and write
\[
 (\pi,\sigma)S=\{(\pi(i),\sigma(j)):(i,j)\in S\}.
\]
Direct substitution into \eqref{eq:Delta} gives the covariance identity
\begin{equation}\label{eq:covariance}
 (\pi,\sigma)\Delta_{f,g}(S)
 =\Delta_{\pi f\pi^{-1},\,\sigma g\sigma^{-1}}
       ((\pi,\sigma)S).
\end{equation}
In particular, renaming coordinates also renames the root. It does
\emph{not} generally give a word from the old root to the renamed subset.

Suppose $R_0\subseteq R$ and $C_0\subseteq C$ contain the root coordinates,
and $S\subseteq R_0\times C_0$. A word on this smaller grid can be lifted
to the large grid by extending each of its functions to fix all excluded
coordinates. Starting in the smaller grid, neither transition branch can
leave it. Hence a reaching word in the smaller grid is also a reaching
word in the full grid.

This removes every empty row or column of a valid target without adding
letters. The distinguished row and column are never empty, so their
removal is not an issue.

\subsection{The one-dimensional base cases}
If $m=1$, a valid subset is $\{0\}\times D$ for a set $D\subseteq C$ that
contains $c$. Starting from $\{(0,c)\}$, add each desired column
$j\in D\setminus\{c\}$ by taking $f=\id$ and letting $g$ send $c$ to $j$
while fixing every other column. The identity branch retains every existing
cell, and the other branch adds only $(0,j)$.

This uses exactly $|S|-1$ letters. The case $n=1$ is symmetric. These are
reachability base cases, not an assertion that the state-complexity formula
is tight for all one-state operands; Section~\ref{sec:limits} explains the
distinction.

\subsection{Deleting a contained copy}
Define the row and column neighborhoods of $S$ by
\[
 N_i(S)=\{j:(i,j)\in S\},\qquad
 M_j(S)=\{i:(i,j)\in S\}.
\]
Suppose no row or column is empty.

\begin{lemma}[Comparable-row reduction]\label{lem:containment}
If $i\ne k$ and $N_i(S)\subseteq N_k(S)$, then, after possibly reversing
$i,k$ in the equality case, there is a valid $T$ with $|T|<|S|$ and a letter
$(f,\id)$ such that $\Delta_{f,\id}(T)=S$. The same statement holds with
rows and columns exchanged.
\end{lemma}
\begin{proof}
Let $(r,c)$ be the root. If $N_i(S)=N_k(S)$ and $k=r$, interchange $i$ and
$k$. After this choice, either $k\ne r$ or the containment is strict. Set
\[
 T=S\setminus\bigl(\{k\}\times N_i(S)\bigr),
 \qquad
 f(x)=\begin{cases}k,&x=i,\\ x,&x\ne i.\end{cases}
\]
The removed set is nonempty. The identity column branch retains $T$.
The row branch copies the retained row $i$ to precisely the removed cells
of row $k$, while producing no cell outside $S$. Thus the image is exactly
$S$.

If $k\ne r$, the distinguished row is unchanged. If $k=r$, strict
containment leaves the nonempty difference $N_k(S)\setminus N_i(S)$ in
that row. Finally, deleting a cell $(k,c)$ cannot empty the distinguished
column: the retained cell $(i,c)$ accompanies every such deletion.
Therefore $T$ remains valid. The column version follows by transposition.
\end{proof}

The equality case is worth retaining in the statement. Blindly deleting
the containing copy can erase the distinguished row when the two rows
are equal. A choice of direction repairs this without changing the
reduction.

Note that the predecessor $T$ of Lemma~\ref{lem:containment} need \emph{not}
have full support: it is a proper subset of $S$ and may have lost a row
entirely. Root-specific validity is enough here precisely because this
reduction is never transported by a relabeling. The contrast with
Definition~\ref{def:certificate}, where full support is demanded, is the
whole reason the two are different objects.

\subsection{A degree-one row or column gives an isolated edge}
After all comparable neighborhoods have been excluded, any degree-one row
or column belongs to an isolated cell. Indeed, if $M_q(S)=\{p\}$ and row
$p$ also meets column $q'$, then $M_q(S)\subseteq M_{q'}(S)$, contrary to
incomparability. The transposed argument handles a degree-one row.

\begin{lemma}[Isolated-cell reduction]\label{lem:isolated}
Assume $m,n\ge2$, all rows and columns of $S$ are nonempty, and $(p,q)$ is
isolated in both its row and its column. A reaching word for
$S\setminus\{(p,q)\}$ in the complementary $(m-1)$-by-$(n-1)$ grid can be
lifted to a reaching word for $S$, with at most two additional letters.
The smaller root can be chosen so that its target is valid.
\end{lemma}
\begin{proof}
Let the original root be $(r,c)$. Choose
\[
 r'=\begin{cases}r,&p\ne r,\\ \text{any element of }R\setminus\{p\},&p=r,
 \end{cases}
 \quad
 c'=\begin{cases}c,&q\ne c,\\ \text{any element of }C\setminus\{q\},&q=c.
 \end{cases}
\]
The target $S_0=S\setminus\{(p,q)\}$ meets every remaining row and column,
so it is valid relative to $(r',c')$.

Let $a$ be the letter whose row function swaps $p,r'$ and whose column
function swaps $q,c'$, fixing all other coordinates. On their two-by-two
rectangle, one application exchanges the two diagonal orientations of a
two-cell set. Starting from the original singleton gives the following
seed words:
\begin{center}
\begin{tabular}{cccc}
\toprule
$p=r$ & $q=c$ & Desired seed & Word \\
\midrule
no & no & $\{(p,q),(r',c')\}$ & $a^2$ \\
yes & no & $\{(p,q),(r',c')\}$ & $a$ \\
no & yes & $\{(p,q),(r',c')\}$ & $a$ \\
yes & yes & $\{(p,q),(r',c')\}$ & $a^2$ \\
\bottomrule
\end{tabular}
\end{center}
These four claims also follow by applying \eqref{eq:cell-transition}
twice explicitly. In the mixed cases the root is on the opposite diagonal
from the desired seed, so one letter, not two, is required.

Take a word reaching $S_0$ from $(r',c')$ in the smaller grid. Extend each
letter to fix $p$ and $q$. The isolated cell then remains fixed, and every
cell of the other component remains within the complementary grid. Since
transitions distribute over union by \eqref{eq:union}, appending this extended
word to the
seed word reaches exactly $S_0\cup\{(p,q)\}=S$.
\end{proof}

\section{Hard matrices and spanning predecessor certificates}
\label{sec:hard}

\subsection{The irreducible targets}
\begin{definition}\label{def:hard}
For $m,n\ge2$, a subset $S\subseteq[m]\times[n]$ is \emph{hard} if every
row and column has at least two cells, distinct row neighborhoods are
incomparable under inclusion, and distinct column neighborhoods are
incomparable under inclusion. Write $\Hard_{m,n}(S)$ for this property.
\end{definition}

The word ``hard'' is local terminology for the remainder after the elementary
reductions, not a computational complexity classification. In particular,
a hard matrix is valid relative to \emph{every} possible root.

A valid target that is not one-dimensional either has an empty row or
column, has comparable neighborhoods, has an isolated degree-one cell
after comparability has been excluded, or is hard. Thus the preceding
lemmas exhaust the non-hard cases.

For a hard target the columns are distinct nonempty proper subsets of
$[m]$. Nonemptiness and distinctness are immediate. A full column would
contain any other column, and $n\ge2$, so a full column is impossible.
Hence, up to column ordering, a hard matrix is a family
\[
 \mathcal A=\{M_0(S),\ldots,M_{n-1}(S)\}
 \subseteq\Pow([m])\setminus\{\varnothing,[m]\}
\]
that is an antichain and also satisfies the row conditions.

\subsection{Why six rows need at most twenty columns}
\begin{lemma}[Antichain width bound]\label{lem:sperner}
An antichain $\mathcal A\subseteq\Pow([m])$ has at most
$\binom{m}{\lfloor m/2\rfloor}$ members. In particular, a hard matrix with
at most six rows has at most twenty columns.
\end{lemma}
\begin{proof}
For a fixed $A\subseteq[m]$, the permutations of $[m]$ whose first $|A|$
entries are exactly the elements of $A$ number $|A|!(m-|A|)!$. Prefix sets
of one permutation form a chain, so an antichain can occur among these
prefix sets at most once. Counting over all $m!$ permutations gives
\[
 \sum_{A\in\mathcal A}|A|!(m-|A|)!\le m!,
 \qquad\text{hence}\qquad
 \sum_{A\in\mathcal A}\binom{m}{|A|}^{-1}\le1.
\]
Each summand is at least $\binom{m}{\lfloor m/2\rfloor}^{-1}$. This proves
the assertion. For $m\le6$ the largest of these middle binomial
coefficients is $\binom63=20$.
\end{proof}

This is the usual antichain bound of Sperner \cite{Sperner1928}, with a
self-contained counting proof included to identify precisely the finite
search space. Applying the same argument to rows also bounds the possible
height for a given width. No assertion about reachability is hidden in the
antichain theorem.

\begin{remark}[The same bound, probabilistically]\label{rem:sperner-prob}
The counting proof above has a one-line probabilistic phrasing: choose a
uniformly random permutation of $[m]$; a fixed $k$-element subset is an
initial segment of it with probability $\binom{m}{k}^{-1}$, and at most one
member of an antichain can be an initial segment, since initial segments are
nested, so $\sum_{A\in\mathcal A}\binom{m}{|A|}^{-1}\le1$. This is the same
mathematics as the proof above, not a second theorem, and is recorded here
only because one of the source computations was documented in that language.
\end{remark}

In fact the finite search space is already bounded without Sperner's theorem
at all. For a hard matrix, every column is a proper subset of $[m]$ of size at
least two, and there are only $2^m-m-2$ such subsets --- fifty-six when $m=6$.
That crude count already makes the recursion over column families finite;
Lemma~\ref{lem:sperner} merely makes it small.

\subsection{A root-independent local proof object}
\begin{definition}[Spanning predecessor]\label{def:certificate}
A spanning predecessor certificate for $S\subseteq[m]\times[n]$ is a triple
$(f,g,T)$ such that
\begin{align}
 &f:[m]\to[m],\qquad g:[n]\to[n],\qquad T\subseteq[m]\times[n],
   \label{eq:certificate-ranges}\\
 &\suppR(T)=[m],\qquad\suppC(T)=[n],\label{eq:certificate-span}\\
 &|T|<|S|,\label{eq:certificate-size}\\
 &\Delta_{f,g}(T)=S.\label{eq:certificate-image}
\end{align}
There is no requirement that $T\subseteq S$.
\end{definition}

If every smaller valid subset is already reachable, such a certificate
proves reachability of $S$ with one additional letter. Full support is
stronger than validity, and this extra strength is what makes the finite
classification economical.

\begin{lemma}[Transport of a certificate]\label{lem:transport}
The existence of a spanning predecessor certificate is invariant under
arbitrary row and column bijections.
\end{lemma}
\begin{proof}
Transport $T$ by $(\pi,\sigma)$ and transport $f,g$ by conjugation.
Cardinality and full support are preserved by bijections; the image
identity is preserved by \eqref{eq:covariance}.
\end{proof}

Notice what the lemma does \emph{not} say: a reaching word from one chosen
root is not asserted to remain a reaching word from that same root after
arbitrary renaming. Instead, the transported \emph{predecessor} is valid
for the actual root, and the induction supplies its reaching word there.
This separates the symmetry argument from the rooted reachability argument.

\subsection{An exact inverse formulation}
For fixed $f,g$ and target $S$, define the allowable source cells
\[
 P_{f,g}(S)=\{(i,j):(f(i),j)\in S\text{ and }(i,g(j))\in S\}.
\]
Every predecessor must be contained in this set. Make a multigraph with
vertex set $S$, placing one tagged edge for each $(i,j)\in P_{f,g}(S)$
between $(f(i),j)$ and $(i,g(j))$. An edge may be a loop, and different
source cells may produce the same pair of endpoints; their source tags
are retained.

\begin{proposition}\label{prop:edge-cover}
A set $T\subseteq P_{f,g}(S)$ satisfies $\Delta_{f,g}(T)=S$ exactly when
the edges tagged by $T$ cover all vertices of this multigraph. The
predecessor is spanning exactly when these edges include every row tag
and every column tag at least once.
\end{proposition}
\begin{proof}
The endpoints of the edge tagged by $(i,j)$ are precisely the two cells
produced by that source. Taking the union of endpoints therefore gives
$\Delta_{f,g}(T)$. The final statement is just the definition of support
in terms of the tags.
\end{proof}

This interpretation explains why local search is much smaller than a
breadth-first search through the entire subset automaton. It also warns
against replacing the problem by an unqualified ordinary edge-cover
problem: the support requirements concern source tags, not only covered
vertices.

The unqualified problem is nevertheless solved exactly, and the solution
makes the warning quantitative.

\begin{proposition}[Minimum unconstrained cover]\label{prop:matching}
Suppose every vertex of the multigraph above is incident with at least one
tagged edge, let $H_0$ be that multigraph with its loops discarded, and let
$\nu(H_0)$ be its maximum matching size. Then the minimum number of tagged
edges covering all vertices is $|S|-\nu(H_0)$.
\end{proposition}
\begin{proof}
For the upper bound, start with a maximum matching and cover each unmatched
vertex by one incident edge or loop. No two unmatched vertices are joined by
a nonloop edge, since that would enlarge the matching; so the number of
selected edges is at most $\nu(H_0)+(|S|-2\nu(H_0))=|S|-\nu(H_0)$.

For the lower bound, take any cover and delete redundant edges. Every
remaining nonloop edge then has an endpoint covered by no other selected
edge, so the selected nonloop components are stars and a selected loop is an
isolated one-vertex component. Choosing one edge from each of the $s$ star
components gives a matching, so $s\le\nu(H_0)$; and the cover has exactly
$|S|-s\ge|S|-\nu(H_0)$ edges.
\end{proof}

Proposition~\ref{prop:matching} is \emph{not} a reachability proof, and the
gap is exactly the support constraint: a minimum edge cover may correspond to
a predecessor missing a row or a column, which is useless for transport by
Lemma~\ref{lem:transport}. The certificates below therefore require full
support explicitly, and the databases of Section~\ref{sec:computation} record
witnesses, not matching numbers.

\subsection{A maximal candidate when the letter is a pair of permutations}
\label{sec:maximal-candidate}
If one restricts attention to letters whose two components are permutations,
the search collapses to arithmetic on finite sets, with no matching solver
and no external solver of any kind. Write $A_j=M_j(S)$ for the target columns.
Then \eqref{eq:Delta} makes the allowable source set of the previous
subsection a column-wise intersection:
\begin{equation}\label{eq:maxpre}
 P_j=f^{-1}(A_j)\cap A_{g(j)},
\end{equation}
and $P$, the subset with these columns, is the \emph{largest} possible
predecessor for the fixed letter $(f,g)$. Its image has columns
\begin{equation}\label{eq:column-image}
 M_j(\Delta_{f,g}(P))=f(P_j)\cup P_{g^{-1}(j)}.
\end{equation}
A candidate is therefore tested with set intersections, permutation images,
and equality of finite sets.

\begin{lemma}[Equal-cardinality repair]\label{lem:repair}
Let $S$ be hard, let $f,g$ be permutations, let $P$ be given by
\eqref{eq:maxpre}, and suppose $\Delta_{f,g}(P)=S$ but $|P|=|S|$. Then some
$P'\subsetneq P$ with $|P'|=|S|-1$ still satisfies $\Delta_{f,g}(P')=S$ and
still has full support.
\end{lemma}
\begin{proof}
Since $P_j\subseteq f^{-1}(A_j)$ for every $j$, the equality $|P|=|S|$ forces
equality in every one of those inclusions; both branches of \eqref{eq:Delta}
are then bijections from $P$ onto $S$. Pick an occupied cell of $P$ whose two
image cells are distinct and delete it. Each of its two images is still
produced, by the other branch applied to a different cell, so the image
remains $S$. Because $S$ is hard, every row and column degree of $S$ is at
least two, and $P$ inherits the same row-degree multiset and the same column
degrees; deleting a single cell therefore cannot empty a row or a column.
\end{proof}

Thus, when the search is restricted to permutation pairs, a strict
predecessor exists whenever the maximal candidate reproduces the target,
including in the tight case. This is the entire discovery procedure used for
database~\dbC{} in Section~\ref{sec:db-residual}, and it needs no solver.
A failure of this restricted strategy on some target would not refute
reachability: a nonpermutation letter, or a route through an intermediate
subset of equal or larger cardinality, might still work. Compare
Section~\ref{sec:opposed}, where a certificate in database~\dbA{} is
exhibited that no permutation-restricted search could find.

\section{A fifth reduction: balanced triples}
\label{sec:triple}
The four reductions so far --- empty axes, containment, isolated cell, and the
certificate step itself --- leave the entire hard class as the finite
obligation. One further \emph{analytic} reduction removes a large part of that
class before any table is consulted. It is not needed for
Theorem~\ref{thm:main}; it gives a second, materially different route to it,
with a much smaller finite remainder.

\begin{definition}[Balanced triple]\label{def:triple}
Three distinct columns of $S$, with neighborhoods $U,V,W\subseteq[m]$, form a
\emph{balanced triple} when
\begin{equation}\label{eq:balanced}
 U\cup V=V\cup W=W\cup U.
\end{equation}
Equivalently: among these three columns, no row occurs exactly once.
\end{definition}

\begin{lemma}[Balanced-triple contraction]\label{lem:triple}
Let $S$ be hard and contain a balanced triple. Then there is a subset $T$ with
full row and column support and $|T|\le|S|-3$, and a letter whose row
component is the identity and whose column component is a three-cycle, with
$\Delta_{\id,g}(T)=S$.
\end{lemma}
\begin{proof}
Let the three columns be $a,b,c$ with neighborhoods $U,V,W$. Define $T$ by
\[
 M_a(T)=U\cap V,\qquad M_b(T)=V\cap W,\qquad M_c(T)=W\cap U,
\]
leaving every other column of $S$ unchanged, and let $g$ be the three-cycle
$a\mapsto b\mapsto c\mapsto a$ fixing all other columns, with $f=\id$.

The output in column $a$ is
\[
 (U\cap V)\cup(W\cap U)=U\cap(V\cup W)=U,
\]
the last equality because $U\subseteq V\cup W$, which is exactly what
\eqref{eq:balanced} says. The other two columns are the same calculation
cyclically, and all remaining columns are fixed by both branches. Hence
$\Delta_{\id,g}(T)=S$.

Each intersection is a proper subset of the corresponding original column,
because the columns are pairwise incomparable; this gives $|T|\le|S|-3$. Each
is also nonempty: if $U\cap V=\varnothing$, then \eqref{eq:balanced} gives
$V\subseteq U\cup W$ and hence $V\subseteq W$, contradicting incomparability.
So no column of $T$ disappears. Finally, a row appearing among these three
columns appears in at least two of them, and any pair of positions in a
three-cycle is one of its consecutive pairs, so that row survives in at least
one of the three intersections; rows occurring only in other columns are
untouched. Therefore $T$ has full support.
\end{proof}

For example, the three columns $\{0,1\},\{1,2\},\{0,2\}$ contract to the three
singleton intersections $\{1\},\{2\},\{0\}$. The identity is
dimension-independent: nothing in the proof mentions six.

\begin{definition}[Residual core]\label{def:residual}
A \emph{residual core} is a hard matrix containing no balanced triple.
\end{definition}

Lemma~\ref{lem:triple} supplies a spanning predecessor certificate, in the
sense of Definition~\ref{def:certificate}, for every hard matrix that is not
residual. Consequently the finite obligation may be narrowed from all hard
matrices to residual cores alone.

\begin{remark}[Hereditary pruning, and what may not be pruned]
\label{rem:hereditary}
Balanced-triple-freeness is \emph{hereditary} under deleting a column: once a
partial column family contains a balanced triple, adding more columns cannot
repair it. This is what allows a column-by-column enumeration of residual
cores to prune a prefix as soon as it acquires a balanced triple, and it is
what makes the resulting search small. Row degree and row incomparability are
emphatically \emph{not} hereditary in this way --- adding a column can repair
a failure of either --- so they are tested at each node and never used to
prune a prefix. This distinction is essential to completeness, and it is the
single place where a plausible-looking optimization would silently break the
enumeration.
\end{remark}

\section{The finite-to-infinite theorem}
\label{sec:induction}
The proof architecture can be stated independently of six and independently
of the method used to find the certificates.

\begin{definition}
For $h\ge1$, let $\mathcal P_h$ be the assertion that every hard matrix with
$2\le m\le h$ rows, of any width, has a spanning predecessor certificate.
By Lemma~\ref{lem:sperner}, this assertion concerns only finitely many
matrices up to coordinate permutations.
\end{definition}

\begin{theorem}[Finite certificate principle]\label{thm:finite-principle}
Assume $\mathcal P_h$. For every $1\le m\le h$, every $n\ge1$, every root
$(r,c)$, and every valid $S\subseteq[m]\times[n]$, there is a word over
$\Gamma_{m,n}$ reaching $S$ from $\{(r,c)\}$. Moreover, it can be chosen
with length at most
\begin{equation}\label{eq:potential}
 B(m,n,S)=|S|+\min(m,n)-2.
\end{equation}
The same conclusion holds when $\min(m,n)\le h$, by transposition.
\end{theorem}
\begin{proof}
We prove reachability and the length bound together, for all roots, by
strong induction on the lexicographically ordered pair
\[
 (m+n,|S|)\in\mathbb N^2.
\]
The induction is well founded. Every dimension-reducing step decreases
its first component; every other recursive step decreases its second.

If $m=1$ or $n=1$, Section~\ref{sec:reductions} gives a word of length
$|S|-1=B(m,n,S)$.

Suppose $m,n\ge2$. If a row or column is empty, delete all empty rows and
columns, keeping the actual root, and obtain a smaller grid of dimensions
$m'\le m$, $n'\le n$, with $m'+n'<m+n$. The cardinality of the target is
unchanged. The induction supplies a word, and its lift adds no letters.
Since $\min(m',n')\le\min(m,n)$, its length is at most $B(m,n,S)$.

We may now assume full support. If there are comparable row or column
neighborhoods, Lemma~\ref{lem:containment} gives a valid predecessor $T$ in
the same grid, with $|T|\le|S|-1$. Reach $T$ by induction and append the
specified letter. The total length is at most
\[
 |T|+\min(m,n)-2+1\le B(m,n,S).
\]

Assume next that all neighborhoods on each side are incomparable. If some
row or column has degree one, the corresponding cell is isolated.
Lemma~\ref{lem:isolated} reduces the problem to a valid target $S_0$ in
an $(m-1)$-by-$(n-1)$ grid, with $|S_0|=|S|-1$. The induction reaches
this smaller target from the chosen smaller root. Adding at most two seed
letters gives a word of length at most
\begin{align*}
 |S_0|+\min(m-1,n-1)-2+2
 &= (|S|-1)+(\min(m,n)-1)\\
 &=B(m,n,S).
\end{align*}

In the remaining case $S$ is hard. Assertion $\mathcal P_h$ gives a
spanning predecessor certificate $(f,g,T)$. Full support of $T$ makes it
valid for the actual root. Because $|T|<|S|$, the induction reaches $T$.
Appending $(f,g)$ reaches $S$, and the same calculation as in the
comparable-neighborhood case proves the length bound.

These cases exhaust every valid target. If only the column dimension is
at most $h$, transpose the matrix, interchange the two function components
of every letter, and apply the result already proved. This preserves both
validity and word length.
\end{proof}

The potential $B$ assigns one unit to every cell and one additional unit
to each available dimension on the smaller side, apart from the base
normalization. An isolated-cell step loses one cell and one unit of the
smaller dimension, paying for its two possible seed letters. This explains
the form of the bound, rather than merely obtaining it from a crude bound
on the number of subset states.

\subsection{The same induction with a weaker finite hypothesis}
\label{sec:criterion-residual}
The proof of Theorem~\ref{thm:finite-principle} consults $\mathcal P_h$ only
in its last case, and only for the target actually in hand. So the hypothesis
may be weakened by anything that discharges part of that case analytically.

\begin{definition}
For $h\ge1$, let $\mathcal R_h$ be the assertion that every \emph{residual}
core (Definition~\ref{def:residual}) with $2\le m\le h$ rows, of any width,
has a spanning predecessor certificate.
\end{definition}

\begin{theorem}[Finite criterion, residual form]\label{thm:finite-residual}
$\mathcal R_h$ implies $\mathcal P_h$, and hence implies every conclusion of
Theorem~\ref{thm:finite-principle}.
\end{theorem}
\begin{proof}
Let $S$ be hard with at most $h$ rows. If $S$ contains a balanced triple,
Lemma~\ref{lem:triple} supplies a spanning predecessor certificate directly.
Otherwise $S$ is residual and $\mathcal R_h$ supplies one.
\end{proof}

This is a strictly weaker finite premise: Section~\ref{sec:db-residual}
discharges $\mathcal R_6$ with 453 certificates over widths at most eleven,
where $\mathcal P_6$ needs more than ten thousand over widths up to twenty.

\subsection{A sharper bound in terms of the support}
\label{sec:support-bound}
\begin{variantbox}{4}
The bound $B(m,n,S)=|S|+\min(m,n)-2$ charges the \emph{grid} dimensions. The
next theorem charges only the dimensions actually occupied. What it buys is
strictly more: on a target with empty axes --- exactly the targets for which
the first reduction of the induction fires --- it is smaller, and it has the
dimension-free corollary $|w|\le2E-2$, which the grid form does not imply. It
is proved by exhibiting an explicit potential and accounting for it case by
case, so it also makes the shape of the bound visible.
\end{variantbox}

\begin{theorem}[Support form of the word bound]\label{thm:support-bound}
Assume $\mathcal P_h$ (or, by Theorem~\ref{thm:finite-residual},
$\mathcal R_h$). Let $1\le m\le h$, let $S$ be valid relative to $(r,c)$, put
$E=|S|$, $r_S=|\suppR(S)|$, and $c_S=|\suppC(S)|$. Then $S$ has a reaching
word of length at most
\begin{equation}\label{eq:support-bound}
 \Phi(S)=E+\min(r_S,c_S)-2 .
\end{equation}
In particular $|w|\le2E-2$, since $r_S,c_S\le E$; and for six rows with
$n\ge6$, $|w|\le6n+4$. The same holds with rows and columns exchanged.
\end{theorem}
\begin{proof}
Run the induction of Theorem~\ref{thm:finite-principle} and account for
$\Phi$ at each step. On a valid singleton $\Phi=0$ and the empty word
suffices. Deleting empty axes changes neither $E$ nor the two support counts,
and adds no letters, so $\Phi$ is unchanged. A containment predecessor is a
proper subset of $S$, so by Corollary~\ref{cor:pred-support} its support
counts do not increase while $E$ drops by at least one; hence
$\Phi(T)\le\Phi(S)-1$, and one letter is appended. A certificate predecessor
has full support, so its support counts equal those of $S$ (which is hard and
hence has full support), and its cardinality is strictly smaller; again
$\Phi(T)\le\Phi(S)-1$ for one appended letter. An isolated-cell step deletes
exactly one cell, one nonempty row, and one nonempty column, so
\[
 \Phi(S')=(E-1)+\min(r_S-1,c_S-1)-2=\Phi(S)-2,
\]
which pays for its at most two seed letters, while the lift adds none. The
one-dimensional base uses $E-1=\Phi(S)$ letters, since one of the support
counts is $1$. Induction on the same rank $(m+n,|S|)$ gives the bound.
\end{proof}

\subsection{A dual bound carried through the induction}
\label{sec:dual-bound}
\begin{variantbox}{5}
The third form carries \emph{two} estimates through the induction at once and
takes their minimum. What the first branch, $2(|S|-1)$, buys is that it does
not mention the dimensions at all, so it is the sharper of the two exactly
when the grid is wide relative to the target; neither branch implies the
other, and the minimum is smaller than either form above on some targets. The
accompanying constructor checks both estimates on every word it emits.
\end{variantbox}

\begin{theorem}[Dual word bound]\label{thm:dual-bound}
Under the hypothesis of Theorem~\ref{thm:finite-principle}, every valid $S$
has a reaching word with
\begin{equation}\label{eq:dual-bound}
 |w|\le\min\bigl\{2(|S|-1),\ |S|+\min(m,n)-2\bigr\}.
\end{equation}
\end{theorem}
\begin{proof}
Write $k=|S|$ and $d=\min(m,n)$, and carry both estimates simultaneously
through the cases of Theorem~\ref{thm:finite-principle}. The singleton uses
zero letters, and the one-dimensional construction uses $k-1$ letters, which
satisfies both. Deleting empty axes changes no cell and adds no letter, while
$d$ cannot increase. Each containment step, balanced-triple step, and
certificate step decreases $k$ by at least one and appends one letter, so both
$2(k-1)$ and $k+d-2$ survive. The isolated-cell step replaces $(k,d)$ by
$(k-1,d-1)$ and adds at most two letters, and its two estimates are
\[
 2\bigl((k-1)-1\bigr)+2=2(k-1),\qquad
 (k-1)+(d-1)-2+2=k+d-2 .
\]
Both bounds therefore hold at every step, hence so does their minimum.
\end{proof}

\subsection{Why this is not an extrapolation from finite experiments}
\label{sec:not-extrapolation}
There is no untested assumption of the form ``the pattern probably continues
for larger $n$.'' Every possible target is classified. A target with empty
axes, comparable supports, or a degree-one axis is reduced analytically; a
target containing a balanced triple is reduced analytically; and every target
not reduced in these ways is one of the finitely many encoded shapes, modulo
row and column relabeling. The finite coverage results then handle all of
those, and the induction applies for any finite number of columns.

Likewise, a certificate table contains predecessors, not an assertion that
those predecessors are already reachable. Their reachability follows from the
strictly decreasing induction. That is why a positive table of local
equations is sufficient to prove a global statement about all reaching words.
The width bounds --- twenty for six rows, or eleven after balanced-triple
removal --- are bounds on the width of an irreducible \emph{obstruction}, not
on the widths the theorem covers. Section~\ref{sec:example} exhibits a
$6\times100$ target handled with no hundred-column certificate.

\section{The finite certificate computations}
\label{sec:computation}
The induction of Section~\ref{sec:induction} needs one finite input:
$\mathcal P_6$, or the weaker $\mathcal R_6$. This section supplies three
separate discharges of that input, from three certificate databases:

\begin{center}
\begin{tabular}{llrrl}
\toprule
& Classified objects & Records & Widths & Discovery\\
\midrule
\dbA{} & hard matrices, orbit-closure coverage & 10,642 & $\le20$ & SMT\\
\dbB{} & hard matrices, degree-sorted canonical & 10,619 & $\le20$ & SMT\\
\dbC{} & residual cores (no balanced triple) & 453 & $\le11$ & solver-free\\
\bottomrule
\end{tabular}
\end{center}

\dbA{} and \dbB{} classify the same objects by two different
symmetry-reduction correctness proofs; \dbC{} classifies a strictly smaller
class, made available by Lemma~\ref{lem:triple}. Any one of the three
suffices for Theorem~\ref{thm:main}. All three share the encoding of the next
subsection.

\begin{variantbox}{0}
What each database buys. \dbA{} is the uncompressed baseline: it certifies the
whole hard class directly, and it is the only one containing certificates
with non-injective row maps, which no permutation-restricted search could
find. \dbB{} is an independent re-derivation of the same finite lemma over a
provably minimal representative set, sharing no witness with \dbA{}, and it
additionally verifies that its row maps are permutations. \dbC{} is the
compressed route: a fifth analytic reduction makes the finite obligation
small enough for a deterministic, solver-free search, so that one path
through the entire result has no external solver anywhere in it.
Section~\ref{sec:db-compare} records precisely why none reduces to another.
\end{variantbox}

\subsection{Encoding the targets}
A column $A\subseteq[m]$ is encoded by
\[
 b(A)=\sum_{i\in A}2^i.
\]
An unordered family of distinct columns is encoded by
\begin{equation}\label{eq:family-mask}
 X(\mathcal A)=\sum_{A\in\mathcal A}2^{b(A)}.
\end{equation}
For $m\le6$, every nonempty proper column has a code between $1$ and $62$.
Thus the entire family fits in an unsigned 64-bit integer. There is no
shift by $64$, and the prohibited full column at code $63$ is not needed.
The family determines a matrix by arranging its column codes in increasing
order. Row permutations act on the binary digits of every column code.

\subsection{Database \dbA{}: orbit-closure coverage}
\label{sec:db-shared}

\subsubsection{The representative list, and why it need not be minimal}
The \dbA{} search generator grows antichains one column at a time and
canonicalizes under row permutations. It first sorts the multiset of row
degrees; among row permutations realizing that degree order, it chooses the
smallest family integer. For completeness of this \emph{search} enumeration,
remove one column from any target antichain: an orbit representative of the
smaller antichain occurs at the preceding level, the same relabeling
transports the removed column to an allowable extension, and canonicalizing
that extension recovers the target orbit. This proves level-by-level coverage
of the generator.

The final proof, however, does not rely on this reasoning or on this
implementation at all. Coverage is audited by an independent checker that
expands the \emph{orbit closure} of each verified representative and compares
against a directly enumerated labeled traversal
(Section~\ref{sec:db-shared-coverage}). Because it compares orbit closures and
the size of the closed set, rather than a list of minimal representatives,
that design is insensitive to redundancy in the representative list: a
representative whose orbit is already covered costs a wasted search but cannot
create a gap, and no minimality of the list is required, or claimed, anywhere
in the argument. This tolerance is a real architectural difference from
database~\dbB{}, whose correctness argument instead \emph{proves} that its
representative list is exactly minimal and whose checker consequently rejects
a duplicate key outright.

A visible consequence of the same difference is that this classification never
transposes. It certifies hard matrices in the orientation in which they arise,
including those with fewer columns than rows, where database~\dbB{} orients
every target with $m\le n$ at the outset and recovers the rest by
transposition. That accounts exactly for the 23 records by which the two
representative sets differ; Section~\ref{sec:db-compare} identifies them.

\subsubsection{The solver query}
For each hard representative, the search introduces integers
$f_0,\ldots,f_{m-1}$ and $g_0,\ldots,g_{n-1}$ with the appropriate finite
ranges, and Boolean variables $t_{ij}$ indicating whether $(i,j)\in T$.
It imposes the following conditions:
\begin{align}
 t_{ij}&\Longrightarrow
       \bigl((f_i,j)\in S\ \land\ (i,g_j)\in S\bigr),
       \label{eq:smt-no-extra}\\
 (r,c)\in S&\Longrightarrow
 \left(\bigvee_{i=0}^{m-1}(t_{ic}\land f_i=r)\right)
 \lor
 \left(\bigvee_{j=0}^{n-1}(t_{rj}\land g_j=c)\right),
       \label{eq:smt-cover}\\
 &\bigvee_jt_{ij}\quad\text{for every }i,\qquad
   \bigvee_it_{ij}\quad\text{for every }j,
       \label{eq:smt-span}\\
 &\sum_{i,j}[t_{ij}]<|S|.
       \label{eq:smt-card}
\end{align}
Membership tests involving $f_i$ or $g_j$ are expanded into disjunctions
of equalities to the constant coordinates of target cells. Here $[t]$ is
$1$ when $t$ is true and $0$ otherwise.

Equation~\eqref{eq:smt-no-extra} forbids extraneous image cells, while
\eqref{eq:smt-cover} requires every target cell to be produced. The last
two lines enforce spanning and strict size reduction. Thus a satisfying
assignment is exactly a spanning predecessor certificate. No optimization
of $|T|$ is required. In the largest shape there are 120 membership
Booleans and 26 function-value integers.

The implementation uses Z3 through its documented C API
\cite{Z3API}, called from Python's standard-library \code{ctypes} module.
The solver is used only to discover witnesses. Its search heuristics,
internal reasoning, and status string are not accepted as the verification
of a certificate. An unsuccessful or timed-out search would also not
refute the shuffle conjecture.

\subsubsection{What was found}
The search returned a spanning predecessor for every one of the 10,642
hard representatives generated for $2\le m\le6$. Of these, 10,573 have six
rows. Table~\ref{tab:coverage} summarizes the independent coverage audit.
A ``labeled family'' here retains row labels but forgets column ordering;
it is not a count of all labeled rectangular matrices.

\begin{table}[htbp]
\centering
\caption{Exhaustive coverage by row count. The antichain column excludes the
special singleton families $\{\varnothing\}$ and $\{[m]\}$, neither of
which can be hard. The empty family is included.}
\label{tab:coverage}
\begin{tabular}{rrrr}
\toprule
Rows & Antichains visited & Hard labeled families & Certificates \\
\midrule
\inputrows tables/coverage_rows.tex 
\bottomrule
\end{tabular}
\end{table}

All target sizes in the certificate file lie between $6$ and $60$;
predecessor sizes lie between $5$ and $38$. Every certificate strictly
reduces size, and the observed reductions range from $1$ to $24$ cells.
These are properties of the concrete file, not unproved universal size
estimates for arbitrary numbers of rows.

\begin{table}[htbp]
\centering
\caption{The six-row hard families, classified by number of columns.
The first numeric column counts row-permutation representatives.}
\label{tab:six-shapes}
\begin{tabular}{rrr}
\toprule
Columns & Certificates & Hard labeled families \\
\midrule
\inputrows tables/six_rows.tex 
\midrule
Total & 10,573 & 5,237,599 \\
\bottomrule
\end{tabular}
\end{table}

The cases below six rows are included rather than treated as black-box
computational results from the literature. This makes the inductive
certificate argument self-contained once the local and coverage audits
are accepted.

The full distribution of six-row hard families by number of columns, for
every $n$ from $0$ through $20$, is given in Appendix~\ref{app:counts}; that
table also gives the total labeled antichain count at each width, which is the
denominator the coverage traversal actually visits. The corresponding
five-row histogram is
\[
 \{4\!:\!30,\ 5\!:\!495,\ 6\!:\!850,\ 7\!:\!445,\ 8\!:\!115,\ 9\!:\!20,\
   10\!:\!2\},
\]
summing to the $1{,}957$ of Table~\ref{tab:coverage}; the four-row histogram
is $\{4\!:\!8,\,5\!:\!6,\,6\!:\!1\}$. All three histograms are recorded in
\code{audit/coverage.json}.

\subsubsection{A solver-independent local checker}
The Python checker reads each JSON record and validates both its structure
and its mathematics. It checks the dimensions, function ranges, distinct
column codes, antichain conditions, row and column degrees, and the family
encoding. It rejects repeated cells in $T$. It then computes the two
images in \eqref{eq:Delta} explicitly and checks
\[
 \suppR(T)=[m],\quad\suppC(T)=[n],\quad |T|<|S|,
 \quad\Delta_{f,g}(T)=S.
\]
All these checks use exact integers and finite sets. Explicit exceptions,
not disableable Python assertions, enforce the independent checker's
conditions. The solver module is not imported by this checker.

After validating a record, the checker exports its target identifier to
\path{audit/verified_targets.tsv}. This ties the subsequent coverage run
to targets whose actual certificates were checked. A list of target names
without validated witnesses would not suffice.

\subsubsection{An independent exhaustive coverage checker}
\label{sec:db-shared-coverage}
The C++ coverage checker does not use the search generator's canonicalizer,
degree ordering, or level-by-level orbit construction. It proceeds in two
separate stages.

First, it expands every verified representative under \emph{all} $m!$ row
permutations, sorts the resulting family integers, and removes duplicates.
Call the resulting orbit-closed set $\mathcal O_m$.

Second, it enumerates every antichain of the nonempty proper subsets of
$[m]$ directly, with row labels fixed. Candidate column codes are kept in
increasing order. At a recursion node, each candidate is considered as the
next column; all candidates comparable to it are deleted before descent.
Every current family is inspected, not just maximal antichains. The
pseudocode is:
\begin{lstlisting}[language=Python]
def visit(family, candidates):
    inspect(family)
    while candidates:
        s = smallest(candidates)
        candidates.remove(s)
        visit(family | {s},
              {t for t in candidates if incomparable(s, t)})
\end{lstlisting}

\begin{lemma}[Coverage enumeration is exhaustive]\label{lem:enumeration}
Starting with the empty family and all nonempty proper column codes, this
recursion visits each antichain exactly once.
\end{lemma}
\begin{proof}
Every recursively added column is larger than the preceding choices and
incomparable with them, so every visited family is an antichain. Conversely,
list the members of an arbitrary antichain in increasing order. At each
step of the recursion, its next member remains a candidate: it is larger
than the members already chosen and is incomparable with all of them.
Following these choices visits the given family. Strictly increasing order
makes this sequence of choices unique.
\end{proof}

\begin{remark}
Two of the source packages phrase this completeness lemma differently --- one
in terms of a set of ``available'' masks from which all comparable masks are
removed on each descent, the other in terms of a candidate set intersected
with the incomparable masks --- but the two phrasings describe the same
recursion and the same proof. Only this one is given. The same recursion,
with one extra pruning rule, is used for database~\dbC{} in
Section~\ref{sec:db-residual}.
\end{remark}

For every visited hard family, the checker demands membership in
$\mathcal O_m$. It also compares the number of hard visited families to
$|\mathcal O_m|$. Thus a missing orbit is not concealed by the fact that
a canonicalization algorithm never generated it. In the six-row run,
all 5,237,599 hard labeled families were in the orbit-closed checked set.
The output is recorded in \code{audit/coverage.json}.

\begin{lemma}[Finite certificate lemma from \dbA{}]\label{lem:P6}
Assertion $\mathcal P_6$ holds.
\end{lemma}
\begin{proof}[Computer-assisted proof]
The independent local checker validated all 10,642 spanning predecessor
certificates in \code{data/certificates.jsonl}. The independent exhaustive
coverage checker, whose enumeration is justified by
Lemma~\ref{lem:enumeration}, verified that every hard column family for
each $m\in\{2,3,4,5,6\}$ is a row permutation of a checked representative.
Column ordering is immaterial by Lemma~\ref{lem:transport}. Every hard
matrix therefore inherits a spanning predecessor certificate. The complete
run reports and executable source are included in the archive.
\end{proof}

\subsubsection{Why the two stages are separated, and a one-binary alternative}
\label{sec:two-stage}
The two-stage design is deliberate. The first stage alone is not a coverage
proof: it validates whatever witnesses it is handed and says nothing about
what is missing. The second stage alone is not a witness proof: it audits a
list of target identifiers, and a list of names without validated witnesses
establishes nothing. The coupling between the stages is the exported file
\path{audit/verified_targets.tsv}, which contains exactly the identifiers of
records that survived stage one, and which stage two consumes as its only
input. Both stages must exit successfully, and the article's finite-to-infinite
argument is a third, separate obligation.

\begin{variantbox}{6}
A second checker architecture for the same database is also shipped:
\code{src/verify_all.cpp}, a single self-contained C++17 binary that performs
the certificate checks, the orbit expansion, and the labeled-antichain
coverage traversal in one process, taking the plain-integer serialization
\code{data/certificates.txt} as its only argument. What it buys is that there
is no intermediate file to get out of step, and no way to run one half of the
check and report the other: the proof of the finite lemma is one command with
one expected line of output. What the two-stage design buys in exchange is
that each stage can be reimplemented, rerun, and audited on its own, and that
the interface between them is an inspectable artifact. Both are kept.
\end{variantbox}

The expected final line of the single binary is
\begin{lstlisting}
PASS: 7836132 labeled antichains enumerated; 5239572 hard antichains covered.
\end{lstlisting}
These are totals across $m=1,\ldots,6$ inclusive, so they are larger than the
$m\ge2$ totals of Table~\ref{tab:coverage}; the one-row case contributes three
antichains and no hard family. Table~\ref{tab:counts-all} gives the breakdown.

\begin{table}[htbp]
\centering
\caption{The single-binary run, including the one-row case. ``Stored'' counts
positive records; no minimality of the representative list is required, or
claimed, for this proof.}
\label{tab:counts-all}
\begin{tabular}{rrrrr}
\toprule
Rows & Stored & All labeled antichains & Hard antichains & Covered\\
\midrule
1 & 0 & 3 & 0 & 0\\
2 & 0 & 6 & 0 & 0\\
3 & 1 & 20 & 1 & 1\\
4 & 5 & 168 & 15 & 15\\
5 & 63 & 7,581 & 1,957 & 1,957\\
6 & 10,573 & 7,828,354 & 5,237,599 & 5,237,599\\
\midrule
Total & 10,642 & 7,836,132 & 5,239,572 & 5,239,572\\
\bottomrule
\end{tabular}
\end{table}

The antichain totals here exceed those of Table~\ref{tab:coverage} by exactly
two per row count, because this traversal also visits the two degenerate
singleton families $\{\varnothing\}$ and $\{[m]\}$, neither of which can be
hard. The two tables are two accountings of one traversal, not two
computations.

\subsubsection{Two serializations, and a test that exists only because of that}
\label{sec:serializations}
Database~\dbA{} is shipped in three byte-different encodings of the same
10,642 witnesses: \path{data/certificates.jsonl}, whose predecessor field
is an explicit list of coordinate pairs; \path{data/certificates_masks.jsonl},
whose predecessor field is a list of column bit masks; and
\path{data/certificates.txt}, a header-commented plain-integer file with one
whitespace-separated record per line, which is what the single binary reads.
Appendix~\ref{app:schema} gives all three layouts.

\begin{variantbox}{7}
Shipping more than one serialization creates a check that is otherwise
impossible: \code{src/test_artifacts.py} compares the plain-integer file and
the mask JSON \emph{record by record}, requiring
$(m,n,\text{key},f,g,\text{predecessor})$ to agree on every line. What it
buys is a cross-check of the writers, not of the mathematics: a transcription
or column-ordering fault that both checkers would accept, because both would
read the same corrupted file, cannot survive agreement between two files
written by different code paths. During the preparation of this merged
archive the same comparison was also run between the coordinate-pair file and
the mask file, converting each record's coordinate list to column masks; all
10,642 records agreed.
\end{variantbox}

\subsection{Database \dbB{}: a minimal canonical representative set}
\label{sec:db-theorem}
Database~\dbB{} discharges $\mathcal P_6$ a second time, over a different
representative set produced by a different symmetry-reduction argument, with
a different search run, a different file format, and different checkers.
It is not a re-verification of \dbA{}; the two share no witness.

\subsubsection{The degree-sorted canonical convention}
Retain a labeled family only when its row-degree sequence is nondecreasing;
among the row permutations preserving that degree sequence, select the one
minimizing the family integer \eqref{eq:family-mask}.

\begin{lemma}[Canonical representatives]\label{lem:canonical}
Every row-permutation orbit of hard matrices contributes exactly one
representative under this convention.
\end{lemma}
\begin{proof}
Any row-degree sequence can be sorted by a row permutation, so every orbit has
at least one retained labeling. All retained labelings in one orbit have the
same sorted degree sequence, so a permutation between two of them sends each
row to a position of equal degree; that is, it permutes only equal-degree
blocks. Minimizing the family integer over exactly those permutations
therefore yields the same integer for every retained labeling in the orbit.
Distinct orbits cannot produce the same encoded family.
\end{proof}

\begin{variantbox}{8}
Lemma~\ref{lem:canonical} and the orbit-closure argument of
Section~\ref{sec:db-shared-coverage} are two different theorems about two
different objects. Lemma~\ref{lem:canonical} says the representative list is
\emph{exactly minimal}, one per orbit, which permits a checker that compares
representative sets directly and rejects duplicates outright. The
orbit-closure argument says nothing about minimality and instead compares
\emph{closures}, which permits a redundant list. Neither implies the other,
and each is the correctness proof of a different shipped program: a checker
built on one assumption would be unsound applied to the other's data file.
Both are kept.
\end{variantbox}

The enumeration itself is the recursion of Lemma~\ref{lem:enumeration}, run
over all $2^m$ column masks, with the hard filters applied at each node:
$n\ge m$, no column of size zero or one, every row degree at least two, and
row incomparability checked in both directions. Column incomparability is
already guaranteed by the recursion.

\subsubsection{Enumeration results}
Table~\ref{tab:totals} gives the totals; the file \code{data/core\_counts.csv}
gives a four-way breakdown per $(m,n)$ --- all antichains, labeled hard
matrices, degree-sorted survivors, and orbits --- which no other count in this
manuscript supplies.

\begin{table}[htbp]
\centering
\caption{Database \dbB{} enumeration, oriented so that $m\le n$. The
``all antichains'' column includes the two degenerate singleton families,
so it exceeds the corresponding column of Table~\ref{tab:coverage} by two.}
\label{tab:totals}
\begin{tabular}{rrrr}
\toprule
Rows $m$ & All antichains & Width bound & Core orbits\\
\midrule
2 & 6 & 2 & 0\\
3 & 20 & 3 & 1\\
4 & 168 & 6 & 5\\
5 & 7,581 & 10 & 62\\
6 & 7,828,354 & 20 & 10,551\\
\midrule
\multicolumn{3}{r}{Total certificate targets} & 10,619\\
\bottomrule
\end{tabular}
\end{table}

The 10,551 six-row orbits are distributed by width as
\begin{center}
\begin{tabular}{rrrrrr}
\toprule
$n$ & Orbits & $n$ & Orbits & $n$ & Orbits\\
\midrule
6 & 296 & 11 & 1,274 & 16 & 25\\
7 & 1,141 & 12 & 706 & 17 & 9\\
8 & 2,133 & 13 & 349 & 18 & 3\\
9 & 2,416 & 14 & 159 & 19 & 1\\
10 & 1,973 & 15 & 65 & 20 & 1\\
\bottomrule
\end{tabular}
\end{center}
so widths $7$ through $20$ account for 10,255 genuinely new representative
targets; adding the 296 six-by-six targets and 68 targets from smaller
dimensions gives 10,619.

\subsubsection{A verifier that regenerates rather than trusts}
The C++20 program \code{src/verify.cpp} both checks every triple and
\emph{reruns the entire canonical enumeration}, then compares the full set of
target keys --- not the number of records --- with the set of keys appearing
in the certificate file. It rejects trailing tokens, duplicate keys, incorrect
column counts, empty source columns, missing source rows, nondecreasing
cardinalities, incorrect images, and a row map that is not a permutation.
Two checkers sharing one target list would not independently establish
coverage; that is exactly why this one regenerates the list.

A second, independently written Python checker,
\code{tests/verify\_certificates.py}, recomputes all 10,619 transition images
using sets of ordered pairs rather than output bit masks, and checks the hard
conditions directly. For $m\le4$ it additionally enumerates \emph{all
families of subsets}, not the antichain recursion, and obtains the same
canonical sets.

\begin{proposition}[Finite certificate lemma from \dbB{}]\label{prop:finite}
Assertion $\mathcal P_6$ holds. Moreover, the row transformation in every one
of the 10,619 supplied certificates is a permutation, and the decreases
$|S|-|T|$ range from $1$ through $22$.
\end{proposition}
\begin{proof}[Computer-assisted proof]
By Lemma~\ref{lem:sperner} only finitely many widths arise. The C++
enumerator implements the recursion of Lemma~\ref{lem:enumeration}, the hard
filters, and the canonicalization of Lemma~\ref{lem:canonical}, producing
10,619 representatives. The verifier checks each supplied triple by the
explicit image computation and then compares key sets: they are equal, with no
missing, duplicate, or extra target. Every triple has full support, strictly
smaller cardinality, the exact required image, and a permutation row map. The
Python checker reproduces all image computations by a different
representation. Lemma~\ref{lem:transport} extends each representative
certificate to all labelings.
\end{proof}

\begin{variantbox}{9}
The additional clause of Proposition~\ref{prop:finite} --- that $f$ is always
a permutation --- is a genuine extra verified property, and it is
\emph{false} of database~\dbA{}, where 8,117 of the 10,642 row maps are not
injective. It matters because a permutation-only search is far cheaper and
needs no solver, and because it is the six-row instance of the stronger
structural statement discussed in Section~\ref{sec:limits}. The column map is
not always a permutation in \dbB{} either: 10,113 of its 10,619 column maps
fail injectivity.
\end{variantbox}

The role of the two checkers must not be conflated: the Python check of
six-row coverage consults the representative manifest, whereas the
\emph{exhaustive} six-row coverage computation is the C++ one.

\subsection{Database \dbC{}: residual cores and no solver at all}
\label{sec:db-residual}
Database~\dbC{} discharges the weaker premise $\mathcal R_6$ of
Theorem~\ref{thm:finite-residual}. Because Lemma~\ref{lem:triple} has already
removed every hard matrix containing a balanced triple, the class that
remains is far smaller, and it is small enough that the discovery procedure of
Section~\ref{sec:maximal-candidate} --- pure permutation enumeration --- is
sufficient.

\subsubsection{The enumeration}
For each $m=2,\ldots,6$, start from the increasing list of masks of proper
subsets of $[m]$ of size at least two. Maintain a chosen family and the
remaining larger candidates. On choosing a new column, reject it if it forms
a balanced triple with two already-chosen columns; otherwise recurse,
retaining only later candidates incomparable with it. At every node with at
least $m$ columns, test row incomparability and row degree, and if they hold
record a residual core.

Completeness is Remark~\ref{rem:hereditary} plus the uniqueness argument of
Lemma~\ref{lem:enumeration}: every residual family has a unique increasing
mask sequence, each of whose prefixes is itself incomparable and
balanced-triple-free, so no prefix can be rejected by either pruning rule and
the family is reached. Row degree and row incomparability are tested at nodes
and never used to prune.

Predecessor columns are stored in the representative's \emph{original} column
order; sorting them independently would generally destroy the certificate.
The records also carry a redundant family key, checked for consistency but
not assumed by the proof.

\begin{table}[htbp]
\centering
\caption{Exact residual-core enumeration, oriented so that $m\le n$. Rows are
labeled and the columns form an unordered family; ``representatives''
additionally quotients by row permutations. Every case with $2\le m\le n$,
$m\le6$, not listed here has no residual core at all.}
\label{tab:cores}
\begin{tabular}{rrrr}
\toprule
Rows $m$ & Columns $n$ & Row-labeled column families & Representatives\\
\midrule
4 & 4 & 7 & 2\\
5 & 5 & 239 & 7\\
5 & 6 & 115 & 4\\
5 & 7 & 10 & 1\\
6 & 6 & 54,358 & 136\\
6 & 7 & 83,730 & 194\\
6 & 8 & 26,565 & 84\\
6 & 9 & 3,355 & 19\\
6 & 10 & 284 & 5\\
6 & 11 & 10 & 1\\
\midrule
\multicolumn{2}{l}{Total} & 168,673 & 453\\
\bottomrule
\end{tabular}
\end{table}

The recursion visits $1,\,7,\,79,\,2{,}715,\,584{,}473$ prefix nodes for
$m=2,\ldots,6$. There are 439 six-row representatives, of which 303 have at
least seven columns; the remaining 14 representatives have fewer than six
rows. The generator tested $690{,}293$ permutation pairs for the six-row
representatives before finding all their certificates.

The enumeration does not stop at eleven columns by assumption: it explores
every admissible family in the whole pool, and the \emph{absence} of residual
cores wider than eleven is an output of the computation, not an input to it.

\subsubsection{A checker that needs no compiler, and a generator that needs no solver}
\begin{variantbox}{10}
\code{residual/verify.py} is a single pure-Python standard-library program
that checks the certificates, expands the row-permutation orbits, checks that
representative orbits are pairwise \emph{disjoint}, and runs the exhaustive
residual enumeration, together with auxiliary checks of the local formulas.
What it buys is that the entire finite lemma can be checked with no C++
compiler, no build step, and no third-party package --- the weakest
prerequisite of any verification path here. Its \code{--quick} flag checks
certificate soundness only and deliberately omits exhaustive coverage; its own
output warns that \code{--quick} must not be represented as verification of
the finite lemma.
\end{variantbox}

\begin{variantbox}{11}
\code{residual/generate.cpp} is a deterministic C++17 regenerator using
integer bit builtins and \emph{no solver of any kind}, implementing the
maximal-candidate formula \eqref{eq:maxpre} with the repair of
Lemma~\ref{lem:repair}. Its regenerated output was compared byte for byte
with the shipped certificate file and found identical. What it buys is that
this one route through the whole result --- discovery and checking --- has
zero dependency on Z3 or on any external solver, so the archive is fully
regenerable as well as fully checkable. The SMT route of
Section~\ref{sec:db-shared} cannot be replaced by it: the shared database
contains non-injective row maps, which a permutation-restricted search cannot
produce.
\end{variantbox}

\begin{lemma}[Finite certificate lemma from \dbC{}]\label{lem:finite-residual}
Assertion $\mathcal R_6$ holds. Moreover, for every residual core with
$2\le m\le6$ the certificate may be taken with \emph{both} $f$ and $g$
permutations.
\end{lemma}
\begin{proof}[Computer-assisted proof]
The file \code{residual/certificates.tsv} contains 453 records, each of which
passes the literal checks of Definition~\ref{def:certificate} together with
bijectivity of $f$ and of $g$. The exhaustive recursion described above covers
every residual core up to column ordering, and every such family lies in
exactly one verified row-permutation orbit; orbit disjointness is checked, and
the number of enumerated families equals the size of the orbit table.
Equation~\eqref{eq:covariance} transports a certificate to any labeling, and
full support and cardinality are invariant under relabeling.
\end{proof}

\subsection{Why none of the three databases reduces to another}
\label{sec:db-compare}
The following comparisons were recomputed directly from the three shipped
files during the preparation of this merged archive.

\begin{enumerate}[label=(\alph*)]
\item \emph{Keys.} \dbB{}'s 10,619 target keys form a strict subset of
\dbA{}'s 10,642. The 23 extra records of \dbA{} are exactly those whose shape
has more rows than columns: twenty-one of shape $6\times5$, one of shape
$5\times4$, and one of shape $6\times4$. \dbA{} certifies each hard matrix in
the orientation in which it arises, whereas \dbB{} orients every target with
$m\le n$ and recovers the others by transposition. Per row count, \dbA{}
stores 1, 5, 63 and 10,573 records at three, four, five and six rows against
\dbB{}'s 1, 5, 62 and 10,551; the differences $1$ and $22$ are the $5\times4$
record and the $6\times5$ and $6\times4$ records respectively. Within each
database every stored record lies in a distinct row-permutation orbit of its
shape, as recomputed for this merge. The numbers 10,642 and 10,619 must not be
``reconciled'' into one: their difference is a real artifact of two different
canonicalization conventions, one of which additionally proves minimality of
its list while the other proves only that its orbit closures cover.
\item \emph{Witnesses.} Of the 10,619 keys present in both databases,
\emph{zero} carry the same triple $(f,g,T)$. Each database is therefore a
genuinely independent re-derivation of the other's finite lemma, not a copy
with a different header.
\item \emph{Structure.} In \dbA{}, 8,117 of the 10,642 row maps are not
injective; in \dbB{}, all 10,619 are permutations while 10,113 column maps are
not; in \dbC{}, all 453 records have both maps bijective.
\item \emph{Class.} \dbC{} classifies a different class of objects. Its
453 records share only two target keys with \dbA{}, and the same two with
\dbB{}.
\item \emph{Remainder.} The two finite classifications of the irreducible
remainder are different objects with different sizes: 5,237,599 six-row hard
labeled families up to width twenty, against 168,673 row-labeled residual
families up to width eleven. Collapsing them would lose either the
balanced-triple compression or the uncompressed baseline that two independent
databases certify.
\end{enumerate}

Consequently the finite input of Section~\ref{sec:induction} is available
three times over, and each way is a complete proof by itself:
Lemma~\ref{lem:P6}, Proposition~\ref{prop:finite}, and
Lemma~\ref{lem:finite-residual} with Theorem~\ref{thm:finite-residual}.
The word constructors and supplementary tests of
Section~\ref{sec:algorithm} provide further checks and usable witnesses, but
are not substituted for any of these coverage arguments.

\section{Completion of the state-complexity proof}
\label{sec:completion}

\begin{proof}[Proof of Theorem~\ref{thm:main}]
Lemma~\ref{lem:P6} and Theorem~\ref{thm:finite-principle} show that every
valid subset is reachable when $\min(m,n)\le6$. Proposition~\ref{prop:finite}
gives the same conclusion through an independent finite computation, and
Lemma~\ref{lem:finite-residual} with Theorem~\ref{thm:finite-residual} gives
it a third time through a smaller one; any one of the three completes this
step. Lemma~\ref{lem:persistence} shows that no invalid subset is reachable.
Hence the universal shuffle automaton has exactly $F(m,n)$ reachable subset
states.

For $m,n\ge2$, give the component automata one final state each.
Section~\ref{sec:distinguishability} proves that the component automata are
minimal with respectively $m,n$ states and that all distinct subset states
are distinguishable --- by any of its four arguments, of which
Corollary~\ref{cor:subsetsdistinct} needs only a fixed three-letter alphabet.
The resulting shuffle language therefore has state complexity $F(m,n)$.
Proposition~\ref{prop:upper} gives the matching upper bound for all operand
pairs. This proves the equality.

Finally,
\begin{align*}
 F(6,n)
 &=2^{6n-1}+31\cdot2^{5n-5}(2^{n-1}-1)\\
 &=63\cdot2^{6n-6}-31\cdot2^{5n-5},
\end{align*}
which proves \eqref{eq:six-simplified}.
\end{proof}

\begin{corollary}[Short reaching words]\label{cor:words}
If $m,n\ge1$ and $\min(m,n)\le6$, every valid subset $S$ of the universal
shuffle state space is reachable by a word of length at most
\[
 \min\bigl\{2(|S|-1),\ |S|+\min(m,n)-2\bigr\},
\]
and also by one of length at most $|S|+\min(r_S,c_S)-2$, where $r_S$ and
$c_S$ are the numbers of nonempty rows and columns of $S$. For a six-by-$n$
grid with $n\ge6$, all three specialize to at most $|S|+4\le6n+4$.
\end{corollary}
\begin{proof}
Apply Theorems~\ref{thm:finite-principle}, \ref{thm:support-bound} and
\ref{thm:dual-bound} with $h=6$.
\end{proof}

Equivalently, the maximum distance from the initial subset to any reachable
subset is at most $mn+\min(m,n)-2$. This is a statement about the universal
alphabet, whose letters already specify complete transformations; it does not
transfer to a fixed generating alphabet, because replacing a transformation
pair by a word in generators need not preserve the subset action.

For scale, the theorem gives
\[
 F(6,6)=66,605,547,520,\qquad
 F(6,7)=4,296,041,037,824,
\]
and already at $n=10$,
\[
 F(6,10)=1,133,816,390,562,611,200 .
\]
The proof did not enumerate any of these reachable-state sets. It
classified the much smaller collection of irreducible targets and used
induction to cover the rest.

\section{Certificates that can be checked by hand}
\label{sec:example}
Four worked six-by-seven examples follow, on \emph{three} different targets,
plus wider examples. They are not illustrations of one phenomenon: the first
two are two incompatible certificates for the same target, drawn from the two
independent databases, and are the cleanest single demonstration that those
databases are independent.

\subsection{A six-by-seven target, certified from database \dbA{}}
\label{sec:ex-shared}
Consider the hard target whose ordered column codes are
\[
 (12,17,18,24,33,34,36).
\]
Its family integer is $94,506,455,040$. With rows and columns numbered
from zero, its incidence matrix and a predecessor are
\begingroup\setlength{\arraycolsep}{3pt}
\[
 S=\begin{pmatrix}
 0&1&0&0&1&0&0\\
 0&0&1&0&0&1&0\\
 1&0&0&0&0&0&1\\
 1&0&0&1&0&0&0\\
 0&1&1&1&0&0&0\\
 0&0&0&0&1&1&1
 \end{pmatrix},\qquad
 T=\begin{pmatrix}
 0&0&0&0&1&0&0\\
 1&0&0&1&0&0&0\\
 0&1&1&0&0&0&0\\
 0&1&0&0&0&0&0\\
 1&0&0&0&0&0&1\\
 0&0&0&0&1&1&0
 \end{pmatrix}.
\]
\endgroup
There are $14$ cells in $S$ and $10$ in $T$. The row and column supports of
$T$ are full. Define the image lists
\[
 f=(0,3,4,0,2,5),\qquad g=(2,0,6,5,4,6,3).
\]
Here the $i$th entry of the first list is $f(i)$, and similarly for $g$.
The two images of each source cell are displayed in
Table~\ref{tab:example-images}.

\begin{table}[H]
\centering
\caption{Direct verification of the six-by-seven predecessor. The union of
the last two columns is exactly the target $S$. Repeated outputs are harmless.}
\label{tab:example-images}
\begin{tabular}{ccc}
\toprule
Source $(i,j)$ & $(f(i),j)$ & $(i,g(j))$ \\
\midrule
\inputrows tables/example_images.tex 
\bottomrule
\end{tabular}
\end{table}

Every output lies in $S$, and every $1$ of $S$ occurs among the outputs.
This checks the image identity without trusting a solver or software.
The predecessor is intentionally not constrained to lie inside its target:
for instance, $(1,0)\in T$ but $(1,0)\notin S$.

\subsection{The same target, certified from database \dbB{}}
\label{sec:ex-theorem}
Database~\dbB{} contains a certificate for this same target --- family integer
$94{,}506{,}455{,}040$ --- and it is a different one. Its row map is a
permutation:
\[
 f=(0,3,5,1,4,2),\qquad
 g=(4,2,2,2,1,0,6),\qquad
 t=(32,16,16,18,1,12,36),
\]
where $t$ lists the predecessor's column codes. As an incidence matrix the
predecessor is
\begingroup\setlength{\arraycolsep}{3pt}
\[
 T'=\begin{pmatrix}
 0&0&0&0&1&0&0\\
 0&0&0&1&0&0&0\\
 0&0&0&0&0&1&1\\
 0&0&0&0&0&1&0\\
 0&1&1&1&0&0&0\\
 1&0&0&0&0&0&1
 \end{pmatrix},
\]
\endgroup
again with ten cells against the target's fourteen, and again with full row
and column support. Separating the two branches of \eqref{eq:Delta} by
column gives
\begin{center}
\begin{tabular}{lrrrrrrr}
\toprule
Column $j$ & 0&1&2&3&4&5&6\\
\midrule
Vertical branch $(f,\id)(T')$ &4&16&16&24&1&34&36\\
Horizontal branch $(\id,g)(T')$ &12&1&18&0&32&0&36\\
Bitwise union &12&17&18&24&33&34&36\\
\bottomrule
\end{tabular}
\end{center}
so the image is exactly $S$. In the \code{.cert} serialization of
Appendix~\ref{app:schema} the whole record is one line:
\begin{lstlisting}
6 7 94506455040 0 3 5 1 4 2 4 2 2 2 1 0 6 32 16 16 18 1 12 36
\end{lstlisting}

\begin{variantbox}{12}
The two certificates just displayed are for the same target and are
incompatible: $f=(0,3,4,0,2,5)$ is not injective, while $f=(0,3,5,1,4,2)$ is a
permutation, and the two predecessors are different ten-cell sets. Neither
can be obtained from the other by the transport of Lemma~\ref{lem:transport},
which acts by conjugation and cannot change whether a map is injective. This
is the concrete form of the fact recorded in Section~\ref{sec:db-compare}:
across all 10,619 shared keys, the two databases agree on no witness at all.
\end{variantbox}

\subsection{An explicit reaching word}
The recursive constructor produces the 11-letter word in
Table~\ref{tab:example-word}, starting at $(0,0)$. The final letter is the
certificate just displayed. The table also gives the number of occupied
cells after each prefix, permitting a simple replay check. The target is
valid even though its initial cell $(0,0)$ is absent.

\begin{table}[H]
\centering
\small
\caption{A reaching word for the six-by-seven example. Image lists specify
the two functions completely. Prefix cardinalities are obtained by exact replay.}
\label{tab:example-word}
\begin{tabular}{rccr}
\toprule
Step & Row function & Column function & Cells \\
\midrule
\inputrows tables/example_word.tex 
\bottomrule
\end{tabular}
\end{table}

The general length bound for this target is $14+6-2=18$, so this word is
comfortably within the proved bound. It is not claimed to be a shortest
word. The same target was also constructed and replayed from every one of
its $6\cdot7=42$ possible roots.

Constructing from the \dbB{} certificate instead produces a different word
for the same target, of ten letters rather than eleven; it is tabulated in
Appendix~\ref{app:theorem-word} and stored in
\path{data/theorem_example_6x7_word.json}. Two constructors reading two
databases give two words; neither length is a claim about optimality.

\subsection{A different six-by-seven target, and a certificate no permutation
search could find}
\label{sec:ex-strip}
Consider now the target whose ordered column codes are
\begin{equation}\label{eq:strip-columns}
 (14,20,25,26,33,34,36),
\end{equation}
so that, for instance, $14=2^1+2^2+2^3$ puts column $0$ in rows $1,2,3$.
It has seventeen cells, meets all six rows and all seven columns, and is
hard; its family integer is $94{,}591{,}008{,}768$. Its certificate in
database~\dbA{} is
\[
 f=(3,5,0,2,4,5),\qquad
 g=(2,2,3,2,0,5,3),
\]
\[
 \bigl(b(M_0(T)),\ldots,b(M_6(T))\bigr)=(9,24,16,17,6,34,10),
\]
with thirteen cells and full support. The two incidence matrices are
\begingroup\setlength{\arraycolsep}{3pt}
\[
 S=\begin{pmatrix}
 0&0&1&0&1&0&0\\
 1&0&0&1&0&1&0\\
 1&1&0&0&0&0&1\\
 1&0&1&1&0&0&0\\
 0&1&1&1&0&0&0\\
 0&0&0&0&1&1&1
 \end{pmatrix},\qquad
 T=\begin{pmatrix}
 1&0&0&1&0&0&0\\
 0&0&0&0&1&1&1\\
 0&0&0&0&1&0&0\\
 1&1&0&0&0&0&1\\
 0&1&1&1&0&0&0\\
 0&0&0&0&0&1&0
 \end{pmatrix}.
\]
\endgroup

\begin{variantbox}{13}
This example is audited in a third style, by columns rather than by source
cells: write $F_j=f(M_j(T))$ for the contribution of the row branch and
$G_j=\bigcup_{k:g(k)=j}M_k(T)$ for the contribution of the column branch.
What that buys is an audit whose row count is the number of \emph{columns},
seven, rather than the number of occupied source cells, thirteen --- a
cheaper hand check on wide targets, and one that displays which branch is
responsible for each column.
\end{variantbox}

\begin{table}[H]
\centering
\caption{Column-wise forward-image check for the target
\eqref{eq:strip-columns}.}
\label{tab:strip-image}
\begin{tabular}{rrrr}
\toprule
$j$ & $b(F_j)$ & $b(G_j)$ & Union, $b(M_j(S))$\\
\midrule
0 & 12 & 6 & 14 \\
1 & 20 & 0 & 20 \\
2 & 16 & 25 & 25 \\
3 & 24 & 26 & 26 \\
4 & 33 & 0 & 33 \\
5 & 32 & 34 & 34 \\
6 & 36 & 0 & 36 \\
\bottomrule
\end{tabular}
\end{table}

From the root $(0,0)$ the constructor produces a sixteen-letter word, printed
in full in Appendix~\ref{app:strip-word} and stored with every intermediate
subset in \code{data/strip\_example\_6x7\_word.json}. Its successive subset
cardinalities are
\[
 1,2,2,3,4,5,5,6,6,7,8,9,10,11,12,13,17 .
\]
The bound here is $17+6-2=21$.

The last of those steps is the certificate letter, and it adds \emph{four}
cells at once using noninvertible transformations. A method restricted to
permutation letters, or to predecessors contained in the target, would miss
this particular certificate. That observation does not claim that every
certificate for this target has the property.

\subsection{A residual core, certified by a bare pair of transpositions}
\label{sec:ex-residual}
The third target is
\begin{equation}\label{eq:residual-columns}
 (3,5,9,14,17,33,50),
\end{equation}
whose columns are $\{0,1\},\{0,2\},\{0,3\},\{1,2,3\},\{0,4\},\{0,5\},
\{1,4,5\}$: sixteen cells, all six rows used, every degree at least two,
both neighborhood families incomparable, and --- checked directly --- no
balanced triple. It is a genuine residual core.

Its certificate in database~\dbC{} is
\[
 f=(0,1,2,3,5,4),\qquad
 g=(0,1,2,3,5,4,6),\qquad
 T=(3,5,9,14,32,17,50),
\]
that is, $f$ swaps rows $4,5$ and $g$ swaps columns $4,5$, and the
predecessor has fifteen cells. This is a minimal one-cell decrease, and the
whole verification is two set identities:
\[
 f(\{5\})\cup\{0,4\}=\{4\}\cup\{0,4\}=\{0,4\},
 \qquad
 f(\{0,4\})\cup\{5\}=\{0,5\}\cup\{5\}=\{0,5\},
\]
giving target columns $4$ and $5$; every other column is fixed pointwise or
setwise, the last one $\{1,4,5\}$ being invariant under the row swap.

\begin{variantbox}{14}
Three worked certificates, three shapes of witness: a non-injective
ten-cell predecessor dropping four cells at once
(Section~\ref{sec:ex-strip}), a permutation-row-map ten-cell predecessor
(Section~\ref{sec:ex-theorem}), and a pair of transpositions dropping one
cell. What the last buys is that it is checkable by hand in two lines with no
table at all, which the other two are not; what the first buys is a witness
outside the reach of any permutation-restricted search.
\end{variantbox}

The constructor of Section~\ref{sec:algorithm} reaches this target from the
origin $(0,0)$ in seventeen letters, tabulated in
\code{residual/example\_6x7.json}; several of its two-letter isolated-edge
seeds preserve cardinality on their second letter while changing positions,
which is why the cardinality trace is not strictly increasing. The word is printed in
Table~\ref{tab:word}.

\begin{table}[htbp]
\centering
\begin{tabular}{rllr}
\toprule
Step & Row transformation $f$ & Column transformation $g$ & Subset size\\
\midrule
1 & \texttt{512340} & \texttt{4123056} & 2\\
2 & \texttt{512340} & \texttt{4123056} & 2\\
3 & \texttt{102345} & \texttt{1023456} & 3\\
4 & \texttt{102345} & \texttt{1023456} & 3\\
5 & \texttt{021345} & \texttt{0321456} & 4\\
6 & \texttt{013245} & \texttt{0213456} & 5\\
7 & \texttt{013245} & \texttt{0213456} & 5\\
8 & \texttt{012435} & \texttt{0153426} & 6\\
9 & \texttt{012435} & \texttt{0153426} & 6\\
10 & \texttt{012345} & \texttt{0126456} & 7\\
11 & \texttt{012345} & \texttt{0133456} & 8\\
12 & \texttt{012345} & \texttt{0323456} & 9\\
13 & \texttt{012345} & \texttt{2123456} & 10\\
14 & \texttt{012345} & \texttt{1123456} & 11\\
15 & \texttt{012345} & \texttt{5123460} & 14\\
16 & \texttt{012345} & \texttt{0123656} & 15\\
17 & \texttt{012354} & \texttt{0123546} & 16\\
\bottomrule
\end{tabular}
\caption{The seventeen-letter reaching word for the residual core, replayed
step by step. The construction is not asserted to be shortest. Several
two-letter isolated-edge seeds preserve cardinality on their second letter
while changing positions.}
\label{tab:word}
\end{table}


\subsection{Two opposed remarks, and their resolution}
\label{sec:opposed}
Two statements in this manuscript look contradictory and are both true.

\begin{enumerate}[label=(\roman*)]
\item Section~\ref{sec:ex-strip}: a method restricted to permutation letters,
or to predecessors contained in $S$, would miss the certificate displayed
there for the hard target \eqref{eq:strip-columns}.
\item Lemma~\ref{lem:finite-residual}: \emph{every} residual core with at most
six rows has a certificate whose two components are both permutations.
\end{enumerate}

The resolution is that the two statements quantify over different targets.
Balanced-triple contraction (Lemma~\ref{lem:triple}) changes which targets
survive to the finite table. The target \eqref{eq:strip-columns} contains a
balanced triple --- its columns $0,1,3$, with neighborhoods
$U=\{1,2,3\}$, $V=\{2,4\}$, $W=\{1,3,4\}$, all three of whose pairwise
unions equal $\{1,2,3,4\}$ --- so on the
\dbC{} route it never reaches the table at all; it is discharged
analytically by a letter whose column component is a three-cycle. What
statement (ii) asserts is a property of the residual remainder, not of every
hard matrix. By contrast the target of
Sections~\ref{sec:ex-shared}--\ref{sec:ex-theorem} has no balanced triple at
all, so it is itself residual; its row-permutation orbit, of size $180$, is
covered in database~\dbC{} by the representative with family key
$282{,}574{,}489{,}649{,}704$, which therefore supplies a third and again
different certificate for it, this one a pair of permutations.
Conversely, database~\dbA{} contains no balanced-triple step,
so it must certify \eqref{eq:strip-columns} directly, and the witness it
found is non-injective. Neither statement may be dropped in favor of the
other: (i) is what forces the SMT route to exist, and (ii) is what makes the
solver-free route possible.

\subsection{Twenty columns and a width beyond the finite classification}
The unique six-row hard family with twenty columns consists of all
three-element subsets of $[6]$. It has $60$ cells. The bundled constructor
reaches the corresponding ordered matrix in $35$ letters; the constructor
driven by database~\dbB{} reaches the same target in $39$ letters. Both are
within the proved bound of $64$. Two constructors, two databases, two words:
the difference measures implementation choices, not the theorem.

Repeat each of these twenty column types five times, obtaining a six-by-one-
hundred matrix with $300$ cells. Repeated columns are handled by the
comparable-column reduction, so no hundred-column certificate is needed.
The supplied reaching word has $115$ letters, compared with the bound
$304$. Both the target and the complete word are included as JSON files.
This example makes the finite-to-infinite step concrete: the finite
classification is of obstruction types, not of every possible target width.

\section{Constructive algorithm and independent tests}
\label{sec:algorithm}

\subsection{Turning the proof into a word}
The supplied \code{construct.py} follows the induction literally. It
accepts a target, dimensions, and a root. It checks validity before starting.
Empty rows and columns are removed with explicit coordinate maps. A
comparable-neighborhood step recursively reaches its smaller valid
predecessor and appends the copying letter. An isolated-cell step prepends
its one- or two-letter seed to a lifted smaller word.

For a hard target, the program tries the at most $6!$ row permutations,
transforms its column codes, and locates a stored representative. The
resulting column permutation is explicit because the column codes are
distinct. If $\pi,\sigma$ carry the current target to a representative
with certificate $(f_0,g_0,T_0)$, the program uses
\[
 T=(\pi^{-1},\sigma^{-1})T_0,\qquad
 f=\pi^{-1}f_0\pi,\qquad g=\sigma^{-1}g_0\sigma.
\]
It checks the transported identity before recursion. On returning, it
replays the entire word from the requested root and checks both the target
and the length bound.

The recursion is mathematically well founded, but this reference Python
implementation uses the interpreter's ordinary recursion stack. Extremely
large user-supplied grids can exceed that implementation limit even though
the theorem applies. No claim of an optimized production implementation is
made.

\subsection{Four constructors with different recursion discipline}
\label{sec:constructors}
Four reaching-word constructors are shipped. They implement the same
induction and are not interchangeable in their interfaces or their checks.

\begin{variantbox}{15}
\code{src/construct\_iterative.py} replaces the recursion by an
\emph{explicit stack} of reduction frames. What it buys is that its logical
recursion depth does not impose the interpreter's recursion limit on the
number of columns, so it is the engine to use on wide targets; it is the only
such implementation here. Running time and memory still grow with the target,
because the transformation arrays are written out explicitly.
\end{variantbox}

\begin{variantbox}{16}
\code{python/construct\_word.py} exposes an importable \code{Constructor}
class with a \code{construct(m,n,S,start=(0,0))} method and a \code{--start
row,column} option, and it reads database~\dbB{}. What it buys is
programmatic reuse: the constructor can be driven from another program
without a JSON round trip, and the root can be varied from the command line.
\end{variantbox}

\begin{variantbox}{17}
\code{residual/construct.py} takes \code{--rows}, \code{--columns},
\code{--masks}, \code{--origin} and \code{--output}, accepts zero masks for
empty columns, raises Python's recursion limit dynamically from the target
size, and checks \emph{both} length bounds of \eqref{eq:dual-bound} on every
word it emits. What it buys is the only implementation that exercises the
balanced-triple step and the only one that validates the dual bound.
\end{variantbox}

\begin{variantbox}{18}
\code{src/construct.py}, the constructor used in
Section~\ref{sec:ex-shared}, takes a target as a JSON object with fields
\code{m}, \code{n}, an optional \code{root}, and \code{target} as a list of
coordinate pairs; it pre-validates the target before starting, then replays
the constructed word and rechecks the length bound before writing output.
What it buys is that the input and the output are both inspectable artifacts,
which is exactly what the bundled examples of Section~\ref{sec:example} are.
The iterative constructor uses the same JSON shape with the key \code{cells}
in place of \code{target}.
\end{variantbox}

All four reject invalid targets, and all four replay the word they produce
and compare the replayed result with the requested target before reporting
success.

\subsection{Exact breadth-first checks on small grids}
A separate bitset breadth-first implementation enumerates the full
transformation-pair alphabet for small grids. It does not call the word
constructor's transition routine. Its reachable-state sets agree exactly
with the independently computed valid-state sets:
\begin{center}
\begin{tabular}{crr}
\toprule
Grid & Reachable subsets & Largest shortest access length \\
\midrule
$2\times2$ & 10 & 2\\
$2\times3$ & 44 & 3\\
$3\times3$ & 400 & 4\\
\bottomrule
\end{tabular}
\end{center}
These breadth-first results check small instances directly. They are
corroboration, not the reason that arbitrarily large widths are covered.

\subsection{Replay, coverage, and negative tests}
The test suite constructed and replayed every valid target on grids
$1\times4$, $4\times1$, $2\times2$, $2\times3$, $3\times3$, and
$3\times4$. The last of these contains $3,392$ valid targets. It also
checked all $3,600$ combinations of a valid three-by-three target and a
specified root, all $42$ roots of the displayed six-by-seven example,
and $180$ reproducibly seeded random targets on larger grids, including
$6\times50$ and its transpose.

An additional straightforward Python antichain enumeration, using lists
and set inclusion rather than the C++ bitset traversal, reproduced both
total and hard-family counts for $m=2,3,4,5$.

Seven deliberately malformed certificates were rejected by the local
checker. The mutations included a truncated function, an out-of-range
function value, an empty predecessor, a repeated source cell, an incorrect
family integer, an unsuccessful solver status, and reversed column order.
A separate negative coverage test removed the sole three-row certificate.
The exhaustive checker then stopped with
\begin{lstlisting}
FAIL: UNCOVERED family: m=3 mask=104
\end{lstlisting}
A test suite that only accepts good data can miss an accidentally vacuous
check. These rejection tests check that important failure paths are active.
They do not replace code review or a formal soundness proof.

\subsection{Three further regression suites}
\label{sec:more-tests}
Each of the other three routes ships its own suite. The four are not
duplicates: they cover different families exhaustively, and their totals are
not comparable.

\paragraph{The iterative constructor's suite: 111,984 replayed words.}
Appendix~\ref{app:tests} gives the per-family breakdown. It exhausts
$1\times1$, $1\times5$, $2\times2$, $2\times3$, $2\times4$, $2\times5$,
$3\times3$ and $3\times4$ over all valid targets \emph{and} all roots ---
the exhaustive $2\times5$ and $4\times4$ families appear in no other suite ---
adds all valid $4\times4$ targets at the root $(0,0)$, relabeled hard targets
with varied roots, random six-row targets up to 64 columns, and the worked
example at all of its roots. A further $2{,}120$ single-cell separation
checks exercise Lemma~\ref{lem:celltest} directly. The same runs record a
reduction-usage census --- $4{,}294$ table reductions, $244{,}114$ row copies,
$130{,}136$ column copies, $219{,}159$ isolated-cell steps, and $176{,}748$
empty-axis removals --- which are implementation diagnostics, not additional
hypotheses.

\paragraph{The \dbB{} suite: 738 access words and 700 distinguishing words.}
The access-word checks exhaust $1\times1$, $1\times4$, $2\times3$,
$2\times4$ and $3\times3$ from the usual root, all valid $2\times2$ targets
from all four roots, and 61 larger cases including randomized targets,
relabeled cores with varied roots, a transposed grid, widths beyond twenty,
and the maximal-width core, with seed $20260920$. The 700
singleton-distinguishing-word tests are the only test suite anywhere here for
the fixed three-letter construction of Section~\ref{sec:three-letter}. Four
named rejection scenarios were exercised against \code{src/verify.cpp}, with
their exact failure strings: a missing record gives \code{exhaustive coverage
failed at m=6}; a duplicated record gives \code{line 10622: duplicate
target}; a negative integer gives \code{only nonnegative decimal integers are
allowed}; and a corrupted witness gives \code{transition image is not the
target}.

\paragraph{The \dbC{} suite: 23,584 replays plus auxiliary local checks.}
It comprises $4{,}201$ tests of all valid origins and targets through
$3\times3$; $19{,}283$ tests of \emph{every certificate target at every
origin}, together with transposes; and 100 deterministic pseudorandom wide
targets at widths $7,11,20,37,64$, twenty samples each, seed $20260920$. The
longest word observed was 133 letters, and both branches of
\eqref{eq:dual-bound} were checked on every word. The verifier additionally
checked $1{,}050$ balanced triples against Lemma~\ref{lem:triple}, $8{,}100$
isolated-edge seeds covering all four origin cases in grids through
$6\times6$, and $8{,}100$ singleton instances of the two-letter cell test.

\paragraph{Encoding and coverage negatives outside the base suite.}
The plain-integer and mask serializations of database~\dbA{} were compared
record by record (Section~\ref{sec:serializations}). Altering one row-map
value to another in-range value was rejected by the single binary because the
forward image no longer matched; deleting the unique three-row certificate
was rejected as the uncovered key $104$ --- the same negative as the base
suite's, reached through a different checker.

\section{Reproduction and the trust boundary}
\label{sec:reproduction}

\subsection{What the archive contains}
The main document is \code{shuffle_six_state.tex} with its compiled PDF, and
\code{tables} holds the TeX tables generated from checked data, so everything
needed to rebuild the PDF is present.

\begin{center}
\footnotesize
\begin{tabular}{@{}>{\ttfamily}l l@{}}
\toprule
\textrm{Path} & Contents\\
\midrule
data/certificates.jsonl & database \dbA{}, coordinate-pair encoding\\
data/certificates\_masks.jsonl & database \dbA{}, column-mask encoding\\
data/certificates.txt & database \dbA{}, plain-integer encoding\\
certificates/cores.cert & database \dbB{}, 10,619 records\\
residual/certificates.tsv & database \dbC{}, 453 records\\
data/cores.csv, core\_counts.csv & \dbB{} target keys and four-way counts\\
src/ & \dbA{} checkers, coverage, constructors, tests, Z3 interface\\
src/verify.cpp, core\_enumeration.hpp & \dbB{} verifier and enumeration header\\
src/enumerate\_cores.cpp & standalone \dbB{} CSV generator\\
python/, tests/, search/ & \dbB{} constructor, checks, optional finder\\
residual/ & \dbC{}: checker, generator, constructor, tests, logs\\
discovery/ & packaged parallel regeneration driver for \dbA{}\\
examples/ & bundled targets, certificates, replayable words\\
audit/, audit/theorem/, logs/ & recorded runs, environments, timings\\
SOURCE\_STATUS.md & sources, attribution and novelty limitations\\
\bottomrule
\end{tabular}
\end{center}

\begin{samepage}
The three certificate files have SHA-256 digests
\begin{center}\small\ttfamily
c44f2ebe73ac2584f1ff9d31cec7fc2e2560ab8fd166c80e86d80eb012c6b594\\
e376820a4de82e1388ea67727ee96bb0cc1c7047d38dd4ae9cfae5617c45abd4\\
6059c807de8aadbadfbdb363c805f83109395063fbbb6e5d87cad45c8b80a48f
\end{center}
\end{samepage}
for databases \dbA{}, \dbB{} and \dbC{} respectively. Each is one
64-character string. A digest identifies exact data; it is not a substitute
for checking them, and this archive deliberately ships no checksum manifest
of its own.

\subsection{Verification without a solver}
The checking workflow requires Python~3.10 or later and a GCC- or
Clang-compatible C++17 compiler. It does not require Z3, a Python package
installation, Internet access, or a proof assistant. From the archive root,
run:
\begin{lstlisting}[language=bash]
make check
\end{lstlisting}
The equivalent core commands are:
\begin{lstlisting}[language=bash]
python3 src/check_certificates.py
c++ -O3 -std=c++17 src/coverage.cpp -o audit/coverage
./audit/coverage audit/verified_targets.tsv > audit/coverage.json
python3 src/test_suite.py
python3 src/negative_coverage_test.py
\end{lstlisting}
Each stage must exit successfully. In particular, a successful local
certificate check alone does not establish coverage. Conversely, exhaustive
coverage of identifiers without checking the associated witnesses does not
establish the predecessor lemma.

A specific example can be regenerated using
\begin{lstlisting}[language=bash]
python3 src/construct.py examples/six_by_seven_input.json \
  --output audit/replayed_example.json
\end{lstlisting}
The supplied PDF can be rebuilt with \code{make pdf}, using a standard
TeX installation with the packages listed in its preamble.

\subsection{Three further verification paths}
\label{sec:paths}
Each database can be checked on its own, and the three routes have
deliberately different prerequisites. None of them needs Z3, a network
connection, a Python package installation, or a proof assistant.

\paragraph{One C++ binary, database \dbA{}.}
\begin{lstlisting}[language=bash]
g++ -O3 -std=c++17 -Wall -Wextra -pedantic src/verify_all.cpp -o verify_all
./verify_all data/certificates.txt
python3 src/test_artifacts.py ./verify_all
python3 src/test_construct_iterative.py
\end{lstlisting}
The final substantive line must be the \code{PASS} line of
Section~\ref{sec:two-stage}. Run the regression script \emph{without}
Python's \code{-O} flag: some of its additional test conditions use the
ordinary assertion mechanism and would be removed. The constructor's own
replay and certificate validation use explicit checks and are not affected
by optimization flags.

\paragraph{One C++ binary, database \dbB{}.}
\begin{lstlisting}[language=bash]
mkdir -p build
g++ -std=c++20 -O2 -Wall -Wextra -pedantic src/verify.cpp -o build/verify
./build/verify certificates/cores.cert
python3 tests/verify_certificates.py
python3 tests/test_construction.py
\end{lstlisting}
The first command both validates the supplied triples and regenerates the
entire canonical representative set, comparing key sets. Expected output ends
with \code{OVERALL: PASS}. The enumeration tables can be regenerated
independently of the certificate file with
\begin{lstlisting}[language=bash]
g++ -std=c++20 -O2 src/enumerate_cores.cpp -o build/enumerate_cores
./build/enumerate_cores build/cores.csv build/core_counts.csv
\end{lstlisting}
and the two CSV files compared with the shipped ones.

\paragraph{No compiler at all, database \dbC{}.}
\begin{lstlisting}[language=bash]
cd residual
python3 verify.py --json verification.json
python3 test_construct.py
\end{lstlisting}
Python 3.9 or newer and its standard library suffice, and the expected final
message is
\begin{lstlisting}
PASS: certificates, exhaustive coverage, and auxiliary checks.
\end{lstlisting}
The \code{--quick} flag checks certificate soundness only, omits
exhaustive coverage, and must not be represented as verification of the
finite lemma. The certificate file can also be regenerated deterministically,
with no solver:
\begin{lstlisting}[language=bash]
g++ -O3 -std=c++17 generate.cpp -o generate
./generate > regenerated.tsv 2> regenerated.log
python3 verify.py --certificates regenerated.tsv
\end{lstlisting}
In the recorded run that regeneration was byte-identical to the shipped file.

\paragraph{Operational notes.}
The command-line constructors accept at most six \emph{rows}. For a grid with
at most six columns instead, transpose the target, construct the word, and
then exchange the two coordinate transformations in every letter; there is no
automatic transposition option. A regeneration run must be given a fresh
output directory, and the \dbA{} search refuses to overwrite an existing
search-output file. Loading the \dbC{} orbit table costs roughly a hundred
megabytes of memory in this environment. On Windows, invoke the Python
entry points through the PowerShell launcher rather than a shell shebang.

\subsection{Optional regeneration of the search}
To regenerate target representatives and search for fresh witnesses, run
\begin{lstlisting}[language=bash]
make enumerate
python3 src/search_certificates.py
\end{lstlisting}
The search requires an installed Z3 shared library. It is located
using \code{ctypes.util.find_library}; an explicit location can be supplied
through the \code{Z3_LIBRARY} environment variable. The default regenerated
certificate output is in \code{audit}, not over the bundled proof data,
and the command refuses to overwrite an existing search-output file.
Different solver versions may choose different witnesses. Such witnesses
are acceptable precisely when they pass the same independent checks.

A ten-target regeneration smoke test was executed with the packaged search
interface, and the packaged target generator exactly reproduced all
10,642 target identifiers. The full original witness-discovery run also
succeeded, as recorded in the audit environment report.

A second, packaged parallel regeneration driver for the same database is in
\code{discovery/}:
\begin{lstlisting}[language=bash]
python3 discovery/regenerate.py regenerated --jobs 4
./verify_all regenerated/certificates.txt
\end{lstlisting}
Its three-row and four-row smoke test is recorded in
\code{logs/regeneration_smoke.txt}. Different solver versions may find
different valid predecessors; bytewise identity with the supplied database is
neither expected nor required for that route. Contrast
Section~\ref{sec:paths}, where the solver-free generator of database~\dbC{}
\emph{was} byte-identical, because its enumeration is deterministic.

The optional finder for database~\dbB{} is \code{search/predecessor_search.py}
with \code{search/z3local.py}; it encodes the same constraints in SMT-LIB and
supports targeted rediscovery by family integer, for example
\code{--rows 6 --family 94506455040 --timeout-ms 10000}.

\subsection{Recorded environment and timing}
The computations used Python~3.13.5, GCC~14.2.0, and Z3~4.13.3.0 on an
x86-64 Linux environment; the database~\dbC{} run used Python~3.13.5 as well.
Recorded elapsed times, by route:

\begin{center}
\small
\begin{tabular}{@{}llr@{}}
\toprule
Route & Step & Elapsed\\
\midrule
\dbA{}, two-stage & initial full witness search & 224.34 s\\
                  & coverage run, peak RSS 67,412 KiB & 3.55 s\\
                  & supplementary test suite & 2.08 s\\
\dbA{}, one binary & witness search after candidate enumeration & 59.54 s\\
                  & original exhaustive checker run & 3.28 s\\
                  & constructor tests & 9.39 s\\
\dbC{}            & verification, peak RSS about 128 MiB & about 7 s\\
\bottomrule
\end{tabular}
\end{center}

These are illustrative observations from these environments, not guaranteed
runtime bounds, portable benchmarks, or a controlled comparison with previous
implementations. The proof depends on the finite checks completing correctly,
not on their speed. Detailed outputs are in \code{audit/},
\code{audit/theorem/}, \code{logs/} and \code{residual/}.

The file \path{logs/verification_ubsan.txt} records a full verification rerun
under undefined-behavior sanitization. It is the only sanitizer run in this
archive; it checks for undefined behavior in the C++ checker's execution,
which ordinary passing output would not reveal, and it is not a proof of the
checker's correctness.

\subsection{Exactly what is, and is not, trusted}
The mathematical reductions are proved in the text. The finite
certificate lemma is supported by explicit witnesses, a local checker,
and an independently structured exhaustive enumeration. The solver is
outside this verification trust boundary: even an incorrect solver cannot
make a false image equality pass a correct checker.

The Python and C++ checkers themselves, their runtimes, the compiler, and
the hardware remain within the practical trust boundary. The several checkers
serve different roles; none is a proof-assistant kernel. The additional
Python enumeration overlaps with the C++ coverage audit only through five
rows, not through six. No claim of independent human review, peer review,
or fully formal verification is made.

The phrase ``independent checker'' throughout this manuscript means
independent of the search and canonicalization algorithms --- and, across the
three databases, independent of one another's representative sets and
witnesses. It does not mean independent human authorship or independently
maintained third-party software: every program here was written during these
investigations. Separately implemented search and checking reduce correlated
implementation risk; they do not eliminate it. Three databases agreeing is
stronger evidence than one, but they were prepared in the same environment,
and that is a residual common cause.

A fully formal version would prove the soundness and completeness of the
finite checker and enumeration inside a proof assistant and import or
recompute the witness table there. Until that is done, the accurate
classification is a reproducible computer-assisted proof, not a
kernel-checked theorem.

\section{What remains open and how to extend the attack}
\label{sec:limits}

\subsection{The unrestricted conjecture}
The argument establishes the strip $\min(m,n)\le6$ and does not establish
$\mathcal P_h$ for any $h\ge7$. Therefore it does not settle the case in
which both operands have at least seven states. Positive computational
results in random larger matrices, even if numerous, would not justify
that extrapolation.

The obstruction reduction survives unchanged at seven rows, where the
width bound becomes $\binom73=35$. However, even the subfamilies of just
the three-element subsets already form $2^{35}$ different antichains.
This elementary lower bound explains why extending the same labeled
coverage traversal without further compression is a materially larger
computation. It is not an argument that the next case is inaccessible.

Two concrete routes are suggested by the present certificates. One is to
classify certificates by structural templates, reducing whole classes of
antichains symbolically. Another is to use a verified isomorph-free
construction in the \emph{coverage} checker itself, while retaining a
separate completeness proof. Either route must preserve the root-validity
invariant that made the current unrooted classification sound. A third is
to look for further analytic reductions in the style of
Lemma~\ref{lem:triple}: that one lemma alone cut the six-row obligation from
5,237,599 labeled families to 168,673, and nothing in its proof mentions six.

\subsection{Three formulations of the open step}
\label{sec:open-formulations}
The unresolved step of this approach can be stated as a conjecture, as a
question, or as a proved conditional. These are three different objects and
all three are recorded.

\subsubsection{As a conjecture}
\begin{conjecture}[Spanning predecessor conjecture]\label{conj:spanning}
Every hard finite zero-one matrix has a spanning predecessor certificate
in the sense of Definition~\ref{def:certificate}.
\end{conjecture}

The present computation proves this assertion for at most six rows or at
most six columns. Theorem~\ref{thm:finite-principle} shows that the general
auxiliary conjecture would imply the full shuffle state-complexity
conjecture, together with the same access-word bound.

The converse is not established. A reachable hard target might have no
smaller \emph{spanning} predecessor and still be reachable from a smaller
merely valid predecessor, or by a path whose final useful step does not
reduce cardinality. Consequently a counterexample to
Conjecture~\ref{conj:spanning} would refute this particular one-step
strategy, not automatically the original shuffle conjecture. A solver
returning ``unknown'' refutes neither statement.

\subsubsection{As a question, with a combinatorial reformulation}
\begin{variantbox}{19}
Stated as a question rather than a conjecture, the same step invites a
reformulation that the conjecture form does not: fix the letter first. What
that buys is an exact finite-combinatorial restatement of the
\emph{existence} problem for a given pair of maps, which is where a
structural proof would have to start.
\end{variantbox}

\begin{question}\label{q:spanning}
Does every hard target, in arbitrary finite dimensions, admit a strictly
smaller full-support predecessor?
\end{question}

Fix $f$ and $g$ first. A candidate source cell $(i,j)$ is allowed exactly
when both $(f(i),j)$ and $(i,g(j))$ lie in $S$; regard its two images as an
edge on the vertex set $S$, a loop when they coincide. Selecting source cells
is then covering all vertices of $S$ by these edges, subject to the extra
requirement that every row and every column be represented among the
\emph{sources}. This is an exact reformulation for fixed maps, and
Proposition~\ref{prop:matching} solves the unconstrained version of it
exactly. The hard part is choosing the maps so that a full-support cover uses
fewer than $|S|$ source cells. This is an observation about the certificate
equations, not a solved structural theorem.

\subsubsection{As a proved conditional}
\begin{variantbox}{20}
The third formulation is not a conjecture at all but a theorem: the
implication itself, stated in two strengths. What it buys over the conjecture
form is that it is \emph{proved}, so it can be cited; and that it isolates
the permutation-pair case, which is what a solver-free search would actually
test.
\end{variantbox}

\begin{proposition}\label{prop:general}
If every residual core in arbitrary dimensions has a strictly smaller
full-support predecessor under some letter of the full transformation
alphabet, then every valid grid subset in arbitrary dimensions is reachable,
and the bounds of Section~\ref{sec:induction} hold there. If the predecessor
letter can always be chosen to be a pair of permutations, the same conclusion
follows as a special case.
\end{proposition}
\begin{proof}
Repeat the induction of Theorem~\ref{thm:finite-principle}, substituting the
hypothesized statement for $\mathcal R_h$ and invoking
Theorem~\ref{thm:finite-residual}; every other reduction, including
Lemma~\ref{lem:triple}, is already dimension-independent. Then apply the
distinguishability argument of Section~\ref{sec:distinguishability}.
\end{proof}

A failed permutation search in a larger dimension would admit three readings,
which must be distinguished. It could refute the permutation-pair form while
leaving a predecessor under a nonpermutation letter --- exactly the situation
of Section~\ref{sec:opposed}. Even failure of \emph{every} one-step strict
predecessor would not prove nonreachability, since a reaching path could
arrive through an intermediate subset of the same or larger cardinality. A
genuine counterexample to the shuffle conjecture requires proving the absence
of \emph{all} reaching words.

The most concrete continuations are therefore a support-constrained version
of the edge-cover problem above, or a structural classification of
balanced-triple-free double-antichain incidence systems. Both are sharply
stated finite-combinatorial questions rather than searches through a
powerset automaton, and neither is claimed to be resolved beyond six rows.

\subsection{Alphabet size and the one-state exception}
All constructive words use the full alphabet $\Gamma_{m,n}$. The result
is not a small-alphabet witness construction. It also does not imply that
an arbitrary pair of $m$- and $n$-state automata has maximal shuffle
complexity: special transition structures can yield much smaller
languages.

The assumption $m,n\ge2$ in the state-complexity theorem is essential to
the proof of distinguishability. That proof requires nonfinal states in
both components. The one-dimensional reachability construction remains
valid, but different reachable subsets can then be equivalent as language
states. Reachability of every valid subset and attainment of their total
number are separate assertions; they must not be conflated.

It is also not proved here that every valid target has a strictly smaller
one-letter predecessor. The induction deliberately keeps the empty-axis and
isolated-cell reductions separate: an isolated-cell seed may need two letters
and may pass through a subset of equal cardinality. The verified one-letter
property concerns only the hard class, and there only with full-support
predecessors.

\subsection{A connection with dependent proof objects}
The artifacts naturally separate discovery from proof checking. In a
dependent-type presentation, a certificate for a fixed target $S$ can be
represented schematically as
\[
 \Cert(S)=
 \sum_{f:R\to R}\ \sum_{g:C\to C}\ \sum_{T:\Pow(R\times C)}
 \left(\Span(T)\times(|T|<|S|)\times
       (\Delta_{f,g}(T)=S)\right).
\]
Here $\Span(T)$ means full row and column support, and the last components
are proofs rather than unchecked flags. The finite computation supplies
concrete candidates for the first three components. A verified decidable
checker would supply the remaining proof components. Well-founded
recursion on $(m+n,|S|)$ then turns these local proof objects into reaching
words with correctness proofs.

This gives a precise formalization target rather than an unverified code
sketch. The nontrivial formal obligations are the enumeration's
completeness, transport of certificates, root-sensitive reductions, and
termination of the constructor. None is delegated to the mere presence
of a solver-generated file.

Two further presentations of the same interface are worth recording. The
first names the certificate type before any checking has occurred,
\[
 \PreCert(S)=
 \sum_{f:[m]\to[m]}\ \sum_{g:[n]\to[n]}
 \bigl\{T\subseteq[m]\times[n]:
   \Span(T)\wedge|T|<|S|\wedge\Delta_{f,g}(T)=S\bigr\},
\]
so that a search \emph{proposes} the computational components $f,g,T$ and a
checker supplies the finite proof obligations that make the proposal inhabit
the specification. The second makes the obligations explicit as named
components of a tuple,
\[
 (T,f,g;\ h_{\mathrm{size}},h_{\mathrm{row}},h_{\mathrm{col}},
  h_{\mathrm{image}}),
\]
asserting respectively $|T|<|S|$, $\suppR(T)=[m]$, $\suppC(T)=[n]$, and
$\Delta_{f,g}(T)=S$; a permutation certificate carries two further
components proving bijectivity of $f$ and of $g$. A small kernel would then
need two theorems and no more: that the certificate-checking Boolean function
implies those properties, and that the recursive enumeration covers every
irreducible target exactly once --- for database~\dbC{} the latter is the
hereditary-pruning argument of Remark~\ref{rem:hereditary}.

The logical shape of the whole result is
\[
 \text{verified finite coverage}
 \ \Longrightarrow\
 \forall n\ge1\ \forall S\subseteq[6]\times[n]:
 \quad\Valid(S)\Longrightarrow\exists w\ \Delta_w(\{(r,c)\})=S .
\]
The universal quantifier over $n$ is discharged by the mathematical
induction, not by evaluating a finite list of column sizes. Root independence
comes from full support, so no hidden quotient by initial-state-changing
permutations is required. Importing only the count 10,619, or 10,642, or 453
as an axiom would not accomplish the coverage proof, which is the part that
matters.

\subsection{Four tempting shortcuts that would be invalid}
Each of the following would look like a proof of Theorem~\ref{thm:main} and
would not be one; each is addressed separately above.
\begin{enumerate}[label=\arabic*.]
\item Numerical agreement with $F(m,n)$ for finitely many pairs $(m,n)$ says
nothing about all $n$.
\item Counting canonical configurations without preserving the
initial-state requirement does not establish reachability: transport moves
the root (Section~\ref{sec:hard}).
\item A satisfiable solver log without witnesses is a weaker artifact than a
checkable certificate file, and is not accepted here at any point.
\item Reachability without distinguishability is not a minimal-DFA lower
bound (Section~\ref{sec:distinguishability}).
\end{enumerate}

\section{Conclusion}
The selected problem was an exact state-complexity conjecture for the
shuffle of regular languages. The outcome is a computer-assisted proof
for the entire six-state strip, supported by finite tables of local
witnesses and exhaustive independent coverage, rather than by a sample
of successful automata calculations.

The essential compression is mathematical: remove empty dimensions,
contained neighborhoods, and isolated cells; optionally contract balanced
triples; classify the remaining antichains; require predecessors to have full
support; and induct on size. The same proof yields explicit reaching words,
of length at most $|S|+\min(m,n)-2$, at most $|S|+\min(r_S,c_S)-2$ in the
support form, and at most $2(|S|-1)$ by the dual estimate.

What this merged archive adds to any one of its sources is redundancy of the
right kind. The finite premise is discharged three times, over two different
classifications of the irreducible remainder, by two different
canonicalization correctness proofs, from two different discovery engines ---
one of which needs no solver --- and audited by four checkers with different
prerequisites, down to one that needs no compiler. Not one witness is shared
between the two databases that classify the same objects. The unrestricted
conjecture remains beyond the result proved here, while the certificate data,
the auxiliary conjecture, its question form, and its proved conditional form
provide concrete next objects for further study.

\appendix
\section{Certificate schemas and minimal verification cores}
\label{app:schema}
Three serialization families are in active use by shipped checkers. They are
not interchangeable, and each is read by a different program.

\subsection{Coordinate-pair JSON lines, and its plain-integer twin}
Each line of \code{data/certificates.jsonl} is an independent JSON object.
The fields \code{m}, \code{n}, and \code{family} give the dimensions and
family integer. The list \code{columns} contains the strictly increasing
column codes. Lists \code{f} and \code{g} are function image tables.
The list \code{T} consists of distinct pairs of zero-based coordinates.
The field \code{status} records that a positive witness was found; its
text is not itself evidence of correctness.

After structural validation, the mathematical core of a local check can
be written as follows. The executable checker contains the additional
range, antichain, degree, and schema checks described in the text.
\begin{lstlisting}[language=Python]
S = {(i, j)
     for j, column in enumerate(record["columns"])
     for i in range(m)
     if column & (1 << i)}
T = {tuple(cell) for cell in record["T"]}
f, g = record["f"], record["g"]

require(len(T) == len(record["T"]), "duplicate cell")
require({i for i, j in T} == set(range(m)), "missing row")
require({j for i, j in T} == set(range(n)), "missing column")
require(len(T) < len(S), "not a smaller predecessor")
image = {(f[i], j) for i, j in T} | {(i, g[j]) for i, j in T}
require(image == S, "incorrect transition image")
\end{lstlisting}

The amount of arithmetic needed to check an individual witness is tiny:
each source cell produces exactly two image cells. The serious global
obligation is not checking one witness, but showing that no irreducible
isomorphism type was omitted. This is the role of the second checker.

The same records appear in \path{data/certificates_masks.jsonl} with the
predecessor given as a list of column bit masks under the key
\code{predecessor} and the family integer under the key \code{key}; and in
\path{data/certificates.txt}, a header-commented file with one
whitespace-separated record per line,
\begin{lstlisting}
m n family_key f[0] ... f[m-1] g[0] ... g[n-1] P[0] ... P[n-1]
\end{lstlisting}
where $P[j]$ is the mask of the predecessor's $j$th column. This is the
form read by \path{src/verify_all.cpp}, which works entirely in column masks:
for every bit $i$ of $P[j]$ it accumulates
\[
 \mathrm{image}[j]\mathrel{|}=2^{f(i)},\qquad
 \mathrm{image}[g(j)]\mathrel{|}=2^{i},
\]
and then compares the resulting array with the target array, with no search
or symbolic inference. Section~\ref{sec:serializations} describes the
agreement test that exists only because both forms are shipped.

\subsection{Fixed-width \texttt{.cert} records for database \dbB{}}
Each noncomment line of \path{certificates/cores.cert} contains exactly
$3+m+2n$ decimal nonnegative integers,
\begin{lstlisting}
m n family f[0] ... f[m-1] g[0] ... g[n-1] t[0] ... t[n-1]
\end{lstlisting}
where \code{family} is \eqref{eq:family-mask}, whose set-bit positions in
increasing order are the target column codes, and the $t_j$ are predecessor
column masks. Comment lines begin with \code{#}, and the file carries the
magic header line \code{# SHUFFLE-FULL-SUPPORT-1}. For example, the unique
three-row hard matrix has columns $3,5,6$ and family integer $104$. The C++
checker validates input domains before shifting bits or indexing arrays; the
additional permutation check on $f$ is specific to this database.

\subsection{Bar-separated records for database \dbC{}}
Each record of \path{residual/certificates.tsv} has five fields separated by
a literal vertical bar,
\begin{lstlisting}
m n family_key | target_columns | row_permutation | column_permutation | predecessor_columns
\end{lstlisting}
with whitespace-separated integer vectors inside each field. Despite the
\code{.tsv} suffix the separator is a bar, not a tab. The predecessor columns
are listed in the representative's original column order and
\emph{must not} be sorted independently; sorting them would in general
destroy the certificate. Both $f$ and $g$ are permutations in every record.

\section{A review checklist}
The following questions isolate the places where a false proof of this
kind could otherwise appear convincing.
{\interlinepenalty=10000
\begin{enumerate}[label=\arabic*.]
\item Are the component functions total and in range? Does the chosen alphabet
contain every letter used in the proof?
\item Does each predecessor remain valid for the root? Are equal neighborhoods
and mixed cases of the isolated-cell reduction handled correctly?
\item Does the hard-target classification forget only symmetries under
which the \emph{certificate property}, not merely a guessed reachability
property, is invariant?
\item Are both inclusions in the image equality checked, along with strict
cardinality reduction and all support requirements?
\item Does coverage include every antichain, not just maximal antichains
or a list of representatives produced by an unverified canonicalizer?
\item Does the state-complexity conclusion include distinguishability and
minimality of the operand automata, rather than just a reachable-state count?
\item Are the finite computation, the infinite induction, the provisional
novelty assessment, and the absence of formal kernel verification clearly
separated?
\end{enumerate}}
Each of these obligations is addressed explicitly in the proof and the
provided audit workflow. Independent reviewers are encouraged to begin
with these points rather than with the solver's search strategy. In this
merged archive, item~5 should be checked separately for each of the three
databases, since they answer it by three different arguments
(Sections~\ref{sec:db-shared-coverage}, \ref{sec:db-theorem} and
\ref{sec:db-residual}).

\clearpage
\section{Six-row antichains by number of columns}
\label{app:counts}
Table~\ref{tab:columns} is taken from the actual exhaustive checker output,
not from an external enumeration table. Degenerate antichains, including the
empty family and the singleton families, are counted in the middle column
even though they are not hard. The middle column is the denominator the
coverage traversal actually visits; the right-hand column is the set of
families that must be certified.

\begin{longtable}{rrr}
\caption{Complete six-row distribution.}\label{tab:columns}\\
\toprule
Columns $n$ & All labeled antichains & Hard labeled antichains\\
\midrule
\endfirsthead
\toprule
Columns $n$ & All labeled antichains & Hard labeled antichains\\
\midrule
\endhead
0 & 1 & 0 \\
1 & 64 & 0 \\
2 & 1,351 & 0 \\
3 & 14,000 & 0 \\
4 & 82,115 & 30 \\
5 & 304,752 & 5,100 \\
6 & 759,457 & 117,821 \\
7 & 1,308,270 & 564,330 \\
8 & 1,613,250 & 1,113,390 \\
9 & 1,484,230 & 1,264,720 \\
10 & 1,067,771 & 998,987 \\
11 & 635,044 & 618,712 \\
12 & 326,990 & 323,900 \\
13 & 147,440 & 147,020 \\
14 & 57,675 & 57,645 \\
15 & 19,238 & 19,238 \\
16 & 5,325 & 5,325 \\
17 & 1,170 & 1,170 \\
18 & 190 & 190 \\
19 & 20 & 20 \\
20 & 1 & 1 \\
\midrule
Total & 7,828,354 & 5,237,599\\
\bottomrule
\end{longtable}

\clearpage
\section{The ten-letter word from database \dbB{}}
\label{app:theorem-word}
This is the word of Section~\ref{sec:ex-theorem} for the target with column
codes $(12,17,18,24,33,34,36)$, constructed from the database~\dbB{}
certificate. Each tuple gives the images of the states in increasing order,
and the letters act by \eqref{eq:Delta}. The initial subset is $\{(0,0)\}$.
The column $|S_k|$ is a diagnostic only: matching cardinalities alone would
not establish that the final set is correct.

\begin{center}
\small
\begin{tabular}{rllr}
\toprule
$k$ & Row map $f_k$ & Column map $g_k$ & $|S_k|$\\
\midrule
1 & $(1,0,2,3,4,5)$ & $(4,1,2,3,0,5,6)$ & 2\\
2 & $(0,2,1,3,4,5)$ & $(3,1,2,0,4,5,6)$ & 3\\
3 & $(0,1,3,2,4,5)$ & $(6,1,2,3,4,5,0)$ & 4\\
4 & $(0,1,2,4,3,5)$ & $(5,1,2,3,4,0,6)$ & 5\\
5 & $(0,1,2,3,5,4)$ & $(1,0,2,3,4,5,6)$ & 6\\
6 & $(0,1,2,3,4,5)$ & $(0,2,2,3,4,5,6)$ & 7\\
7 & $(0,1,5,3,4,5)$ & $(0,1,2,3,4,5,6)$ & 8\\
8 & $(0,1,2,2,4,5)$ & $(0,1,2,3,4,5,6)$ & 9\\
9 & $(0,4,2,3,4,5)$ & $(0,1,2,3,4,5,6)$ & 10\\
10 & $(0,3,5,1,4,2)$ & $(4,2,2,2,1,0,6)$ & 14\\
\bottomrule
\end{tabular}
\end{center}


The same data in machine-readable form is
\code{data/theorem_example_6x7_word.json}. For an independent direct check,
start from the singleton, apply the two image operations at every step, and
compare the final seven column codes.

\clearpage
\section{The sixteen-letter word for the second six-by-seven target}
\label{app:strip-word}
Each $f$-tuple lists the images of rows $0$ through $5$ and each $g$-tuple
lists the images of columns $0$ through $6$; the letters are read from top to
bottom. The final column gives the number of occupied cells after that
letter, starting from the singleton $\{(0,0)\}$. The exact subsets, not only
their sizes, are stored in \code{data/strip_example_6x7_word.json} and are
replayed by the constructor.

\begin{table}[H]
\centering
\small
\caption{A sixteen-letter reaching word for the target
\eqref{eq:strip-columns}.}
\label{tab:strip-word}
\begin{tabular}{rllr}
\toprule
Step & $f$ & $g$ & Cells\\
\midrule
1 & $(1,0,2,3,4,5)$ & $(1,0,2,3,4,5,6)$ & 2 \\
2 & $(1,0,2,3,4,5)$ & $(1,0,2,3,4,5,6)$ & 2 \\
3 & $(0,2,1,3,4,5)$ & $(0,6,2,3,4,5,1)$ & 3 \\
4 & $(0,1,3,2,4,5)$ & $(0,4,2,3,1,5,6)$ & 4 \\
5 & $(0,1,2,4,3,5)$ & $(0,2,1,3,4,5,6)$ & 5 \\
6 & $(0,1,2,4,3,5)$ & $(0,2,1,3,4,5,6)$ & 5 \\
7 & $(0,1,2,3,5,4)$ & $(0,1,5,3,4,2,6)$ & 6 \\
8 & $(0,1,2,3,5,4)$ & $(0,1,5,3,4,2,6)$ & 6 \\
9 & $(0,1,2,3,4,5)$ & $(0,1,3,3,4,5,6)$ & 7 \\
10 & $(0,1,2,3,4,5)$ & $(3,1,2,3,4,5,6)$ & 8 \\
11 & $(0,1,2,3,4,5)$ & $(0,0,2,3,4,5,6)$ & 9 \\
12 & $(0,1,2,3,4,5)$ & $(0,1,1,3,4,5,6)$ & 10 \\
13 & $(0,3,2,3,4,5)$ & $(0,1,2,3,4,5,6)$ & 11 \\
14 & $(0,1,2,3,4,1)$ & $(0,1,2,3,4,5,6)$ & 12 \\
15 & $(0,1,1,3,4,5)$ & $(0,1,2,3,4,5,6)$ & 13 \\
16 & $(3,5,0,2,4,5)$ & $(2,2,3,2,0,5,3)$ & 17 \\
\bottomrule
\end{tabular}
\end{table}

\clearpage
\section{Detailed regression-test counts}
\label{app:tests}
The following counts describe the executed deterministic suite of
Section~\ref{sec:more-tests}. Tests with different roots are counted
separately, since validity and reaching words depend on the distinguished
axes.
\begin{center}
\begin{tabular}{lr}
\toprule
Test family & Replayed words\\
\midrule
$1\times1$, all valid targets and roots & 1\\
$1\times5$, all valid targets and roots & 80\\
$2\times2$, all valid targets and roots & 40\\
$2\times3$, all valid targets and roots & 264\\
$2\times4$, all valid targets and roots & 1,472\\
$2\times5$, all valid targets and roots & 7,520\\
$3\times3$, all valid targets and roots & 3,600\\
$3\times4$, all valid targets and roots & 40,704\\
$4\times4$, every valid target, root $(0,0)$ & 57,856\\
Relabeled hard targets, varied roots & 240\\
Random six-row targets, up to 64 columns & 75\\
Matching targets and all roots of the worked example & 132\\
\midrule
Total & 111,984\\
\bottomrule
\end{tabular}
\end{center}
The random generator is seeded by the integer $20260920$.

\section{Proof dependency summary}
\label{app:dependency}
The logical flow of the argument is:
\[
\begin{gathered}
\text{Validity invariant}+\text{inclusion--exclusion}\\
\Downarrow\\[-5pt]
\text{Universal upper bound }F(m,n),\\[4pt]
\text{Structural reductions}+\text{antichain bound}
 +\text{complete enumeration}\\
\Downarrow\\[-5pt]
\text{$\mathcal P_6$ by \dbA{}, or by \dbB{};\quad
       $\mathcal R_6$ by \dbC{} with Lemma~\ref{lem:triple}}\\
\Downarrow\\[-5pt]
\text{Reachability for }\min(m,n)\le6
 \text{ and three access-word bounds},\\[4pt]
\text{Reachability}+\text{any one of four distinguishability arguments}\\
\Downarrow\\[-5pt]
\text{Sharp deterministic shuffle state complexity.}
\end{gathered}
\]
The prior results for at most five states and for $6\times6$ position the
work historically; they are not used as unverified computational assumptions
anywhere in this chain.

\begin{thebibliography}{9}
\addcontentsline{toc}{section}{References}
\bibitem{BJLRS2016}
Janusz Brzozowski, Galina Jir\'askov\'a, Bo Liu, Aayush Rajasekaran,
and Marek Szyku\l a.
\newblock On the state complexity of the shuffle of regular languages.
\newblock In \emph{Descriptional Complexity of Formal Systems (DCFS 2016)},
Lecture Notes in Computer Science 9777, pages 73--86. Springer, 2016.
\newblock DOI: \href{https://doi.org/10.1007/978-3-319-41114-9_6}{10.1007/978-3-319-41114-9\_6}.
\newblock Accessible version:
\href{https://arxiv.org/abs/1512.01187v3}{arXiv:1512.01187v3}, 15 July 2016.

\bibitem{CLP2020}
Pascal Caron, Jean-Gabriel Luque, and Bruno Patrou.
\newblock A combinatorial approach for the state complexity of the shuffle product.
\newblock \emph{Journal of Automata, Languages and Combinatorics},
25(4):291--320, 2020.
\newblock DOI: \href{https://doi.org/10.25596/jalc-2020-291}{10.25596/jalc-2020-291}.
\newblock Accessible preprint:
\href{https://arxiv.org/abs/1905.08120}{arXiv:1905.08120}, 20 May 2019.
\newblock The preprint, rather than the access-restricted journal full text,
was used for the mathematical source review.

\bibitem{Sperner1928}
Emanuel Sperner.
\newblock Ein Satz \"uber Untermengen einer endlichen Menge.
\newblock \emph{Mathematische Zeitschrift}, 27:544--548, 1928.
\newblock DOI: \href{https://doi.org/10.1007/BF01171114}{10.1007/BF01171114}.

\bibitem{Z3API}
The Z3 contributors.
\newblock Z3 C API documentation.
\newblock \url{https://z3prover.github.io/api/html/group__capi.html}.
\newblock Consulted 20 September 2026. The computation used the locally
installed shared library reporting version 4.13.3.0.
\end{thebibliography}
\end{document}
